\documentclass[11pt]{article}
\usepackage[margin=1.1in]{geometry}
\usepackage{amsmath, amssymb, amsthm}
\usepackage{graphicx}
\newtheorem{proposition}{Proposition}[section]
\newtheorem{theorem}[proposition]{Theorem}
\newtheorem{corollary}[proposition]{Corollary}
\newcommand{\E}{\mathbb{E}}
\newcommand{\Prob}{\mathbb{P}}
\newcommand{\1}{\mathbf{1}}
\usepackage{booktabs}
\usepackage{array,longtable,pdflscape}
\usepackage{microtype}
\newcommand{\thead}[1]{\shortstack{#1}}
\usepackage{xr-hyper}
\usepackage[hidelinks]{hyperref}
\hypersetup{pdftitle={Supplement to Learned representations and kernel corrections for radiative-transfer emulation},pdfauthor={Yitzchak Shmalo}}
\graphicspath{{../figures/}}
\externaldocument[][nocite]{emit_kernel_dnn}[https://github.com/yspennstate/emit-oco2-kernel-emulators/raw/main/paper/emit_kernel_dnn.pdf]

% Sections, results, tables, figures and equations of this document carry the prefix S.
\renewcommand{\thesection}{S\arabic{section}}
\renewcommand{\thetable}{S\arabic{table}}
\renewcommand{\thefigure}{S\arabic{figure}}
\numberwithin{equation}{section}

\title{Supplement to ``Learned representations and kernel corrections\\ for radiative-transfer emulation''}
\author{Yitzchak Shmalo\thanks{Einstein Institute of Mathematics, The Hebrew University of Jerusalem, Jerusalem, Israel. \texttt{yitzchak.shmalo@gmail.com}.}}
\date{September 2026}

\begin{document}
\maketitle

\begin{abstract}
This supplement contains the proofs of the results stated in the article, three further results on the
sensitivity, selection and perturbation of learned corrections, and the experiments that support the article
without being needed to follow it: further representation and combination studies on the EMIT table, a synthetic
check of the transfer theorem, a matched correction-coefficient benchmark, full-data tuning and gradient-based
metrics, a comparison on OCO-2, ten further emulation configurations and extended tables. Sections, results,
tables and equations with the prefix S are in this document; numbers without it refer to the article.
\end{abstract}

\section{Sensitivity, selection and perturbation of learned corrections}\label{supp:sensitivity}

\subsection{Sensitivity of the correction}\label{sec:sensitivity}

For a common frozen operator $R=I-S_\lambda$ and the norm of Section~\ref{sec:identity}, the exact criterion
$2\langle Rf,Rh\rangle>\|Rh\|^2$ is necessary and sufficient for strict improvement of the corrected predictor over
the standalone kernel, since $\|Rf-Rh\|^2<\|Rf\|^2$ is that inequality rearranged; the comparison of smoothness
norms is neither necessary nor sufficient. Write the design values of the targets, of the fitted mean and of the
corrected predictor as $y$, $h_X$ and $\widehat y$, with $\widehat y=h_X+S_\lambda(y-h_X)$ and
$S_\lambda=K(K+\lambda I)^{-1}$; in the nugget convention of the experiments, $\lambda=n\gamma$. The mean is trained
on the same labels, so $h_X=h_X(y)$. In the following Jacobian calculation,
$K$ and $\lambda$ are held fixed; label-dependent feature or kernel selection would add
terms involving derivatives of $S_\lambda$.

\begin{proposition}\label{prop:sensitivity}
\begin{enumerate}
\item If $h_X$ is differentiable in $y$, the Jacobian of this fixed-kernel pipeline is
$J_{\widehat y}=S_\lambda+(I-S_\lambda)J_{h}$. In particular the correction stage alone does not bound the
sensitivity of the pipeline: with $K=I$, $\lambda=1$ and $h_X(y)=2y$ one has $\widehat y=\tfrac32 y$ although
$\|S_\lambda\|_{\mathrm{op}}=\tfrac12$; and if the mean interpolates the training labels, $h_X(y)=y$, then
$\widehat y=y$ and the correction changes nothing on the design.
\item With the mean and the kernel frozen, $y\mapsto\widehat y$ is affine with linear part $S_\lambda$, its
effective degrees of freedom are $\operatorname{tr}S_\lambda=\sum_j\mu_j/(\mu_j+\lambda)$, and
$\|S_\lambda\|_{\mathrm{op}}=\max_j\mu_j/(\mu_j+\lambda)<1$.
\item Suppose the mean, the kernel and the ridge are fixed independently of the correction noise, conditional on
the design, and the correction is fitted on an independent sample with $y=f_X+\varepsilon$, $\E\varepsilon=0$,
$\E[\varepsilon\varepsilon^{\top}]=\sigma^2I$. With $r=f_X-h_X$ and $r_j$ its coordinates in the
eigenbasis of $K$,
\[
\frac1n\,\E_\varepsilon\|\widehat y-f_X\|^2
=\frac1n\sum_j\Bigl(\frac{\lambda}{\mu_j+\lambda}\Bigr)^2r_j^2
+\frac{\sigma^2}{n}\sum_j\Bigl(\frac{\mu_j}{\mu_j+\lambda}\Bigr)^2 .
\]
\item At a new input $x$, with $a_\lambda(x)=(K+\lambda I)^{-1}k(X_n,x)$, the noise variance of the correction
is $\sigma^2\|a_\lambda(x)\|_2^2$.
\end{enumerate}
\end{proposition}

\noindent\emph{Proof in Section~\ref{app:proofs}.}
Item (iii) is what the ridge-path experiments report: the residual's spectral energy $\sum_j r_j^2$ by eigendirection,
and the two sums as functions of $\lambda$, together with the out-of-sample factor of item (iv). A ridge chosen from
fits that depend on the correction noise is outside the hypothesis of (iii).

\subsection{A finite-class selection guarantee}\label{sec:damped}

A finite comparison on a genuinely independent selection sample has a simple risk bound
\cite{stonecv,gcv,arlot}. The independence assumption is stronger than reusing the same
validation rows for training decisions and combination weights.

\begin{proposition}\label{prop:selection}
Let $p_1,\dots,p_M$ be predictors frozen before an independent selection sample of $m$ i.i.d.\ units, let the loss
take values in $[0,B]$, and let $\widehat p$ minimise the empirical risk on that sample. Then with probability at
least $1-\delta$,
\[
R(\widehat p)\le\min_{1\le j\le M}R(p_j)+2B\sqrt{\frac{\log(2M/\delta)}{2m}} .
\]
\end{proposition}

\noindent\emph{Proof in Section~\ref{app:proofs}.}
The squared constrained-retrieval loss is bounded by $R^2$. Thus the proposition can
compare a finite set of frozen corrections or stacks, with bands and prescribed
reflectances averaged within each independent state. For example, a finite grid of
damped corrections $h+\alpha k(x,X_n)(K+\lambda I)^{-1}(y-h_X)$ can include $\alpha=0$;
the standalone kernel must be added as a separate candidate when it is not already in
the grid. To apply the result to stacks, one first fits their weights on development
rows and then compares a frozen finite collection on an independent state-level sample.
In the experiments of the article the same validation rows serve early stopping, kernel tuning and
combination, so the proposition describes the protocol under which such a comparison is
certified rather than a property of the reported stacks.

\subsection{Perturbing a learned-mean kernel predictor}\label{sec:perturb}

Fix the training sample and write a corrected predictor as
$p(x)=h(x)+k_x^{\top}H^{-1}(Y-M)$ with $H=K+\lambda I$ positive definite, $Y\in\mathbb R^{n\times q}$ the targets
and $M\in\mathbb R^{n\times q}$ the training-row means used as correction labels. Let a second fitted system be
$\widetilde p(x)=\widetilde h(x)+\widetilde k_x^{\top}\widetilde H^{-1}(Y-\widetilde M)$ with $\widetilde H$
symmetric positive definite; both systems may depend on $Y$ through the mean, the features, the selected ridge and
any kernel approximation. Put $B=H^{-1/2}(\widetilde H-H)H^{-1/2}$, whose least eigenvalue exceeds $-1$,
$Q=(I+B)^{-1}$, $\gamma_Q=\|Q\|_{\mathrm{op}}=1/(1+\lambda_{\min}(B))$,
$\eta_Q=\|Q-I\|_{\mathrm{op}}=\max_i|b_i|/(1+b_i)$ over the eigenvalues $b_i$ of $B$, and
\[
g=H^{-1/2}k_x,\quad \widetilde g=H^{-1/2}\widetilde k_x,\quad v=\widetilde g-g,\quad
W=H^{-1/2}(Y-M),\quad U=H^{-1/2}(\widetilde M-M).
\]

A conditioning argument bounds such a difference by the condition number of $H$ times a norm of the
change \cite{stewartsun,golubvanloan,higham}. The statement below instead keeps the residual action
$W$ and the label shift $U$ separate, which is what makes it readable on data.

\begin{proposition}[same-data perturbation]\label{prop:perturb}
\begin{enumerate}
\item The exact difference is
\[
\widetilde p(x)-p(x)=\widetilde h(x)-h(x)+v^{\top}W
                  -(Q\widetilde g)^{\top}(BW+U).
\]
\item $\|\widetilde p(x)-p(x)\|_2\le\|\widetilde h(x)-h(x)\|_2+\|W^{\top}v\|_2
+\|Q\widetilde g\|_2\,\|BW+U\|_{\mathrm{op}}$.
\item $\|\widetilde p(x)-p(x)\|_2\le\|\widetilde h(x)-h(x)\|_2
+\|W\|_F\bigl(\gamma_Q\|v\|_2+\eta_Q\|g\|_2\bigr)+\gamma_Q\|U\|_F\|\widetilde g\|_2$.
\item If $\widetilde H=H+sI$ with $s\ge0$ and the kernel and the means unchanged, then
$\|\widetilde p(x)-p(x)\|_2\le\|g\|_2\|W\|_F\,s/(\lambda_{\min}(H)+s)$.
\end{enumerate}
\end{proposition}

\noindent\emph{Proof in Section~\ref{app:proofs}.}
Only positive definiteness of $\widetilde H$ is needed. A large positive eigenvalue of $B$,
as can arise from increasing the ridge, does not invalidate the bound. The bound can
nonetheless be numerically large, especially when the least eigenvalue approaches $-1$. Item (ii) keeps the
interaction between the perturbation and the residual: with $v=0$, $U=0$ and $BW=0$ the kernel correction does not move, however large $\|B\|$ is; the full predictor is
unchanged as well when $\widetilde h(x)=h(x)$. The bounds are stated in the residual $W$, not in $Y$, they use no independence between
the representation, the mean, the ridge and $Y$, and they transfer between fitted predictors without being
generalisation statements.

For component $j$ the physical target at band $b$ is $f_{j,b}=\mu_{j,b}+\sum_\ell(A_j)_{b\ell}z_{j,\ell}+\pi_{j,b}$,
with $A_j$ the standardiser scale times the retained PCA basis and $\pi_j$ the truncation residual, and the
pipeline predicts the retained coordinates. For a frozen realised system
$p_\ell(x)=h_\ell(x)+a_x^{\top}(y_\ell-h_{\ell,X})$ with $y_\ell=f_{\ell,X}+\varepsilon_\ell$, the conditional
RKHS evaluation bound \cite{rw,nmkc} gives
\begin{equation}\label{eq:chain-coord}
|p_\ell(x)-f_\ell(x)|\le B_\ell\widetilde P_\lambda(x)+|a_x^{\top}\varepsilon_\ell|,
\end{equation}
where, writing $k_x=k(X_n,x)$ and $a_x=(K+\lambda I)^{-1}k_x$,
\[
\widetilde P_\lambda(x)^2
 =k(x,x)-2a_x^\top k_x+a_x^\top K a_x
 =k(x,x)-k_x^\top(K+\lambda I)^{-1}k_x-\lambda\|a_x\|_2^2.
\]
This follows by applying the RKHS Cauchy--Schwarz inequality to the residual and
$k(x,\cdot)-\sum_i(a_x)_i k(x_i,\cdot)$.
Here $B_\ell$ is a bound on the residual RKHS norm, a hypothesis and not an observed certificate, and the second term is zero for
deterministic simulator targets, which is the form used for EMIT. A perturbed system adds the bound
$\Delta_\ell(x)$ of Proposition~\ref{prop:perturb}. If $|\widehat z_{j,\ell}-z_{j,\ell}|\le u_{j,\ell}$, then
\begin{equation}\label{eq:chain-recon}
|\widehat f_{j,b}-f_{j,b}|\le|\pi_{j,b}|+\sum_\ell|(A_j)_{b\ell}|\,u_{j,\ell},
\end{equation}
and a change of basis or standardiser between two systems adds its own reconstruction term. Assembling
$e_R=|e_a|+R|e_t|+R^2t|e_s|$ from \eqref{eq:chain-recon} for the three physical components and applying
Theorem~\ref{thm:transfer} closes the chain. If the inverse is evaluated at $L_{\mathrm{obs}}=L+\eta$, the proof of
Theorem~\ref{thm:transfer} goes through with $e_a$ replaced by $e_a-\eta$, so
$|\bar\rho-\rho|\le\min\{R,(|e_a-\eta|+R|e_t|+R^2t|e_s|)/t\}$. The theorem is an
exact-arithmetic statement. In floating point the added term $\rho t/q$ can be lost when it is
small relative to the spacing of representable numbers near $a$, in a way that depends on
$\rho$, $q$, the magnitudes and the rounding rather than on a single threshold such as
$t<2^{-52}a$; a rounded radiance may then give zero for $L-a$ even with the tabulated
components. A count of numerical violations is therefore a check on the implementation, not a
test of the inequality.



\section{Representation, correction and combination}\label{sec:results-representation}

\paragraph{Concatenated features.}
Table~\ref{tab:v2-replication} gives the three-member width-2000 experiment of
Section~\ref{sec:scaling} split by split, with the networks trained on one GPU.

\begin{table}[p]
\centering
% Generated by code/format_publication_tables.py; numerical cells preserved.
{\fontsize{9}{11}\selectfont
\setlength{\tabcolsep}{2.8pt}
\begin{tabular}{llrrrr}
\toprule
seed & head & radiance [\%] & median [pp] & p95 [pp] & gain [pp] \\
\midrule
104 & feature kernel & 0.0924 & 0.044 & 9.25 & -- \\
104 & concatenated features & 0.0810 & 0.025 & 6.25 & +0.0114 \\
104 & convex stack & 0.0781 & 0.024 & 12.77 & -- \\
104 & network + residual & 0.1843 & 0.115 & 54.71 & -- \\
105 & feature kernel & 0.0611 & 0.031 & 7.31 & -- \\
105 & concatenated features & 0.0508 & 0.028 & 6.02 & +0.0103 \\
105 & convex stack & 0.0475 & 0.024 & 14.50 & -- \\
105 & network + residual & 0.1349 & 0.087 & 42.61 & -- \\
106 & feature kernel & 0.0789 & 0.043 & 7.67 & -- \\
106 & concatenated features & 0.0702 & 0.035 & 6.28 & +0.0087 \\
106 & convex stack & 0.0687 & 0.034 & 16.29 & -- \\
106 & network + residual & 0.1608 & 0.108 & 51.84 & -- \\
107 & feature kernel & 0.0731 & 0.042 & 5.82 & -- \\
107 & concatenated features & 0.0732 & 0.044 & 5.41 & -0.0002 \\
107 & convex stack & 0.0833 & 0.043 & 13.84 & -- \\
107 & network + residual & 0.1476 & 0.103 & 46.58 & -- \\
mean & feature kernel & 0.0764$\pm$0.0130 & 0.040$\pm$0.006 & 7.51$\pm$1.41 & -- \\
mean & concatenated features & 0.0688$\pm$0.0128 & 0.033$\pm$0.009 & 5.99$\pm$0.40 & 0.0076$\pm$0.0053 \\
mean & convex stack & 0.0694$\pm$0.0158 & 0.031$\pm$0.009 & 14.35$\pm$1.48 & -- \\
mean & network + residual & 0.1569$\pm$0.0211 & 0.103$\pm$0.012 & 48.93$\pm$5.40 & -- \\
\bottomrule
\end{tabular}}

\caption{Kernel on concatenated features, three width-2000 networks, splits 104--107. Radiance
is a mean relative spectral error in percent; reflectance median and 95th-percentile errors are
in percentage points. The last column is the paired radiance gain of concatenated features over
the features of one network. The mean row averages the four within-split quantities.}
\label{tab:v2-replication}
\end{table}

The concatenated-feature kernel has mean radiance error $0.0688\%$, against $0.0764\%$ for the
features of one network; the paired gain is $0.0076\pm0.0053$ percentage points, positive on
three splits and $-0.0002$ on the fourth. Its mean within-split 95th-percentile reflectance error
is $5.99$ points, against $7.51$ for single-member features and $14.35$ for the stack. With the
ten width-512 splits of Section~\ref{sec:scaling} this makes two settings, with different
networks on overlapping splits, in which concatenation lowers both the radiance error and the
reflectance tail.

\paragraph{Kernel choice on the state and on learned features.}
Table~\ref{tab:v2-multikernel} covers four components on seeds 101--103. The spectral mixture
and the Mat\'ern-5/2 kernel have similar mean errors on $Y_1$ ($0.089\%$), and the spectral
mixture has the lowest mean errors of the raw-state kernels on $Y_2$ and $Y_3$ ($0.237\%$ and
$0.099\%$). The first-order additive kernel is much less accurate in this specification.
Albedo errors for the nonadditive kernels lie between about $2.3\%$ and $3.5\%$.

\begin{table}[p]
\centering
% Generated by code/format_publication_tables.py; numerical cells preserved.
{\fontsize{9}{11}\selectfont
\setlength{\tabcolsep}{3pt}
\begin{tabular}{lrrrr}
\toprule
head & $Y_1$ [\%] & $Y_2$ [\%] & $Y_3$ [\%] & $Y_4$ [\%] \\
\midrule
Mat\'ern-1/2 & 0.179$\pm$0.008 & 0.672$\pm$0.003 & 0.418$\pm$0.012 & 2.326$\pm$0.103 \\
Mat\'ern-3/2 & 0.092$\pm$0.002 & 0.391$\pm$0.008 & 0.207$\pm$0.016 & 2.481$\pm$0.043 \\
Mat\'ern-5/2 & 0.089$\pm$0.017 & 0.434$\pm$0.007 & 0.203$\pm$0.003 & 2.746$\pm$0.082 \\
Gaussian & 0.092$\pm$0.002 & 0.525$\pm$0.113 & 0.228$\pm$0.051 & 3.502$\pm$0.501 \\
spectral mixture (4) & 0.089$\pm$0.010 & 0.237$\pm$0.014 & 0.099$\pm$0.013 & 3.468$\pm$0.263 \\
NNGP, ReLU & 0.175$\pm$0.001 & 0.674$\pm$0.015 & 0.350$\pm$0.009 & 2.548$\pm$0.057 \\
NTK, ReLU & 0.263$\pm$0.002 & 1.055$\pm$0.015 & 0.621$\pm$0.014 & 2.321$\pm$0.048 \\
additive Mat\'ern-5/2 & 4.091$\pm$0.093 & 20.266$\pm$1.035 & 11.430$\pm$0.237 & 4.413$\pm$0.077 \\
ARD by kernel flow & 0.322$\pm$0.015 & 0.257$\pm$0.019 & 0.199$\pm$0.005 & 2.418$\pm$0.114 \\
ARD by marginal likelihood & 0.281$\pm$0.039 & 0.341$\pm$0.042 & 0.183$\pm$0.011 & 2.702$\pm$0.027 \\
convex kernel sum & 0.087$\pm$0.007 & 0.253$\pm$0.020 & 0.099$\pm$0.013 & 2.557$\pm$0.099 \\
network & 0.819$\pm$0.027 & 2.289$\pm$0.019 & 0.910$\pm$0.057 & 3.482$\pm$0.085 \\
network + residual kernel & 0.523$\pm$0.007 & 1.337$\pm$0.026 & 0.588$\pm$0.035 & 3.055$\pm$0.121 \\
features, Mat\'ern-5/2 & 0.105$\pm$0.006 & 0.211$\pm$0.016 & 0.114$\pm$0.008 & 2.303$\pm$0.216 \\
features, Gaussian & 0.126$\pm$0.007 & 0.262$\pm$0.018 & 0.138$\pm$0.014 & 2.426$\pm$0.214 \\
features, NTK & 0.152$\pm$0.005 & 0.307$\pm$0.015 & 0.156$\pm$0.007 & 2.150$\pm$0.059 \\
features, ARD & 0.116$\pm$0.007 & 0.214$\pm$0.014 & 0.111$\pm$0.004 & 2.217$\pm$0.207 \\
coordinate selection & 0.090$\pm$0.010 & 0.360$\pm$0.008 & 0.105$\pm$0.010 & 2.797$\pm$0.447 \\
convex stack of the heads & 0.066$\pm$0.003 & 0.156$\pm$0.010 & 0.077$\pm$0.004 & 1.956$\pm$0.037 \\
ridge stack of the heads & 0.087$\pm$0.005 & 0.288$\pm$0.031 & 0.087$\pm$0.005 & 2.310$\pm$0.071 \\
\bottomrule
\end{tabular}}

\caption{Ten input-space kernels, feature kernels and combination rules on the EMIT state
inputs, with rank-64 outputs. Entries are component mean relative spectral errors in percent,
averaged over seeds 101--103 with sample standard deviations. The network has width 512 and depth
3 with a 150-epoch budget. Kernels include the Mat\'ern family, Gaussian, spectral mixture
\cite{wilsonadams}, NNGP \cite{chosaul,leenngp}, NTK \cite{ntk}, additive kernel
\cite{additivegp}, and learned input scales \cite{owhadiyoo,rw}.}
\label{tab:v2-multikernel}
\end{table}

The Mat\'ern-5/2 feature kernel gives $0.105\%$, $0.211\%$, $0.114\%$ and $2.303\%$ on
$Y_1$--$Y_4$, better than the best raw-state kernel on some components and not on others. The
convex stack has the lowest component means in this comparison ($0.066\%$, $0.156\%$, $0.077\%$,
$1.956\%$), which is consistent with its large reflectance tails elsewhere: the component metric
does not weight errors the way the inverse does.

\paragraph{Network regularization.}
Table~\ref{tab:v2-regsuite} compares a control, a weight-decay grid and an early-stopping grid on
three seeds. The validation-selected weight decay matches the control's component mean
($2.27\%$), and the selected early stopping is less accurate and more variable
($2.72\pm0.45\%$). Residual correction brings the component means to $1.88$--$1.95\%$ and
feature kernels to $0.95$--$0.99\%$, with radiance errors near $0.15\%$, whichever
regularization the network used. The reference row is the isotropic kernel fitted in the same runs.
Under this budget, regularizing the network changes little; the kernel stage does the work.

\begin{table}[t]
\centering
% Generated by code/format_publication_tables.py; numerical cells preserved.
{\fontsize{9}{11}\selectfont
\setlength{\tabcolsep}{3pt}
\begin{tabular}{lrrr}
\toprule
head & \thead{alone\\component / radiance} & \thead{residual correction\\component / radiance} & \thead{feature kernel\\component / radiance} \\
\midrule
isotropic kernel & \shortstack{0.941$\pm$0.109\\0.284$\pm$0.078} & -- & -- \\
network, control & \shortstack{2.266$\pm$0.031\\0.809$\pm$0.029} & \shortstack{1.903$\pm$0.055\\0.617$\pm$0.033} & \shortstack{0.955$\pm$0.038\\0.153$\pm$0.052} \\
network, weight decay & \shortstack{2.269$\pm$0.027\\0.820$\pm$0.014} & \shortstack{1.880$\pm$0.072\\0.618$\pm$0.040} & \shortstack{0.948$\pm$0.048\\0.154$\pm$0.048} \\
network, early stopping & \shortstack{2.716$\pm$0.445\\1.293$\pm$0.373} & \shortstack{1.949$\pm$0.073\\0.789$\pm$0.127} & \shortstack{0.986$\pm$0.032\\0.147$\pm$0.054} \\
\bottomrule
\end{tabular}}

\caption{Regularization comparison, seeds 101--103. Each cell gives the mean component
relative error [percent] / radiance relative error [percent], with sample standard deviations.
The width-512, depth-3 network has a 150-epoch cap, and validation selects among three runs per
arm. The weight-decay grid is $10^{-5},10^{-4},10^{-3}$ and the early-stopping patience grid 3, 5
and 10 epochs \cite{adamw,prechelt}. The reference row is the isotropic kernel fitted in the same
runs.}
\label{tab:v2-regsuite}
\end{table}

\paragraph{Perturbation and ridge-path diagnostics.}
The tables of Section~\ref{app:diagnostics} give the perturbation diagnostics for the state
inputs and the learned features. The residual-action bound of Proposition~\ref{prop:perturb} is
often much smaller than the spectrum-only bound, yet it can still exceed the actual movement by a
factor of tens, thousands or more: positive definiteness makes the bound applicable without
making it a useful tolerance.

Table~\ref{tab:v2-ridge} reports a 4{,}000-row fit of four components at seed 101. At the
smallest ridge on input-space $Y_1$ the effective degrees of freedom are 3{,}841 and the noise
factor $0.924$. The same-half and split-half principal-component errors are $0.079\%$ and
$2.465\%$, and the physical-unit test errors of the corrected and standalone kernels $0.2168\%$
and $0.1754\%$; these columns use different metrics and samples and are not to be subtracted. At
$\gamma=10^{-2}$ the corresponding test errors are $0.5095\%$ and $5.7578\%$: the fitted mean
keeps its predictive information when the kernel correction is strongly regularized.

\section{A synthetic check of the transfer theorem}\label{sec:results-synthetic}

Table~\ref{tab:v2-synthetic} checks the construction of Proposition~\ref{prop:attained} on
manufactured data, two million samples for each $\beta\in\{0.5,1,2,3,4\}$. The average-error
slopes are $0.200$, $0.333$, $0.498$, $0.597$ and $0.664$, against the limits $0.200$, $0.333$,
$0.500$, $0.600$ and $0.667$. For the uniform family the finite-ladder slope is compared with the
analytic curve on the same ladder: at $\beta=2$ the measured slope is $1.721$ and the analytic one
$1.724$, while the limit is two. The largest relative Monte Carlo discrepancy is about $1.01\%$ for
the uniform family at $\beta=2$ and below $0.053\%$ for the profile family against its quadrature
reference, and no pointwise violation occurs among the ten million sampled transmissions.

\section{The matched benchmark and the transfer contrast}\label{sec:results-benchmarks}

\paragraph{Correction-coefficient benchmark.}
Table~\ref{tab:v2-pkanrtm} scores our families under the standard and out-of-distribution
protocols of \cite{pkan}, with and without the low-fidelity 6S coefficients as inputs, against the
published values of the reference model, which we did not retrain.

The state-only residual-corrected network reaches standard-split RMSE $0.00581\pm0.00004$, MAE
$0.00151$, uniform-average $R^2=0.9952$ and SMAPE $3.54\%$, against published values $0.00619$,
$0.00186$, $0.9947$ and $3.91\%$; with 6S inputs its means are $0.00578$, $0.00148$, $0.9954$ and
$3.19\%$. On the out-of-distribution split the state-only means are $0.00620$, $0.00252$, $0.9967$
and $3.37\%$, against published $0.01022$, $0.00544$, $0.9942$ and $5.40\%$, and with 6S inputs
$0.00596$, $0.00232$, $0.9969$ and $2.92\%$. The spreads describe three seeds on the standard
protocol and three initializations on the out-of-distribution protocol.

The $R^2$ convention matters here. The reference weights the three coefficients equally, whereas a
variance-weighted $R^2$ is the default of many libraries. We report the uniform average, computed
from per-coefficient target scales that are recovered from the per-coefficient RMSE of the heads
and checked against the variance-weighted value. The ensemble without correction is close to the
corrected network, while the raw-state kernel is much less accurate in these configurations.

\paragraph{OCO-2 coefficient and radiance objectives.}
Table~\ref{tab:v2-oco2} covers three bands and seeds 0--2. The feature kernel trained under the
flat coefficient loss has coefficient errors of about $4.1\%$, $16.3\%$ and $8.1\%$ on O$_2$, weak
CO$_2$ and strong CO$_2$, and radiance errors of $0.076\%$, $0.114\%$ and $0.108\%$. A
radiance-weighted feature construction gives about $0.031\%$, $0.048\%$ and $0.064\%$ in radiance,
at the price of larger coefficient errors. The reconstruction map and the evaluation metric decide
which fitted representation is preferred.

The radiance stack minimizes a row-relative squared radiance loss, a convex quadratic because the
reconstruction is affine, and its solver reaches KKT gaps between $4.17\times10^{-6}$ and
$1.52\times10^{-5}$. Its mean radiance errors are $0.0264\%$, $0.0391\%$ and $0.0511\%$, against
$0.0291\%$, $0.0404\%$ and $0.0536\%$ for the squared-loss stack on reduced coordinates, while its
coefficient errors rise to $18.59\%$ and $37.43\%$ on O$_2$ and weak CO$_2$ and to $8.87\%$ on
strong CO$_2$. All rules use 2{,}000 validation and 2{,}000 test rows.


\section{Full-data tuning and gradient-based metrics}\label{sec:retune}

We compare hyperparameter selection on a 6{,}000-row subsample with selection on all
training rows, and coordinate search with a metric read off the gradients of a cubic surrogate
(Table~\ref{tab:retune}).

\begin{table}[t]
\centering\small
\resizebox{\textwidth}{!}{\begin{tabular}{lcccccc}
\toprule
 & \multicolumn{4}{c}{component rel.\ $L^2$ [\%]} & radiance & reflectance \\
\cmidrule(lr){2-5}
kernel & $Y_1$ & $Y_2$ & $Y_3$ & $Y_4$ & rel.\ $L^2$ [\%] & med.\ $|\hat\rho-\rho|$ [\%] \\
\midrule
isotropic, tuned on 6{,}000 rows & 0.08$\pm$0.01 & 0.42$\pm$0.03 & 0.20$\pm$0.00 & 2.62$\pm$0.18 & 0.239$\pm$0.016 & 0.141$\pm$0.018 \\
isotropic, tuned on every row & 0.08$\pm$0.00 & 0.36$\pm$0.02 & 0.17$\pm$0.01 & 2.58$\pm$0.14 & 0.210$\pm$0.008 & 0.107$\pm$0.006 \\
per-input search, tuned on 6{,}000 rows & 0.08$\pm$0.01 & 0.16$\pm$0.02 & 0.08$\pm$0.01 & 2.18$\pm$0.16 & 0.095$\pm$0.009 & 0.036$\pm$0.004 \\
per-input search, re-tuned on every row & 0.08$\pm$0.00 & 0.16$\pm$0.02 & 0.08$\pm$0.01 & 2.09$\pm$0.12 & 0.096$\pm$0.009 & 0.036$\pm$0.004 \\
per-input from the cubic's gradients, $p=1$ & 0.34$\pm$0.01 & 0.19$\pm$0.02 & 0.14$\pm$0.01 & 2.17$\pm$0.14 & 0.159$\pm$0.007 & 0.078$\pm$0.004 \\
per-input from the cubic's gradients, $p=1/2$ & 0.11$\pm$0.00 & 0.19$\pm$0.02 & 0.09$\pm$0.01 & 2.27$\pm$0.15 & 0.125$\pm$0.008 & 0.055$\pm$0.004 \\
per-input from the cubic's gradients, $p=1/4$ & 0.08$\pm$0.00 & 0.25$\pm$0.02 & 0.14$\pm$0.01 & 2.38$\pm$0.20 & 0.157$\pm$0.015 & 0.077$\pm$0.013 \\
\bottomrule
\end{tabular}
}
\caption{Where the kernel's hyperparameters are chosen, and how its metric is obtained, at the
full training block and the same ten splits. The first pair tunes smoothness, scale and nugget
on a $6{,}000$-row subsample and on every training row; the last three rows replace the coordinate
search by a metric read off the cubic surrogate's gradients, at three exponents. Selecting the
exponent on validation, per split, chooses $p=1/2$ on all ten, so that row would repeat the
$p=1/2$ row and is omitted. Mean $\pm$ standard deviation over the ten splits.}
\label{tab:retune}
\end{table}

\paragraph{Tuning with every row in the solve.} The kernel families of
Section~\ref{sec:methods} choose smoothness, scale and nugget by the validation error of a solve
on a $6{,}000$-row subsample and refit the winner on all $18{,}884$ rows. Tuning with the full
solve instead lowers the isotropic kernel's radiance error from $0.239\%$ to $0.210\%$ and its
median reflectance error from $0.141$ to $0.107$ points, better on nine of ten splits, with the
ratio averaging $0.882$ and never above $1.001$. The selected scales differ: over the forty
component-split cells the full block picks a shorter length scale than the subsample in $26$,
the same in $12$ and a longer one in $2$. A smaller tuning fit prefers more smoothing.

Below $6{,}000$ rows the tuning subsample is the whole training block, so those rungs are
already fully tuned; at $8{,}000$ rows it is three quarters of the block and only at the top rung
is it a third. Applying the measured ratio to the top rung of Table~\ref{tab:curves} takes the
kernel's last step on radiance from $-0.183$ to $-0.334$, against a slope of $-0.305$ over the
five lower rungs, and on the two flux components from $-0.187$ and $-0.171$ to $-0.352$ and
$-0.358$, against $-0.437$ and $-0.468$. Most of the bend of Section~\ref{sec:scaling} is the
tuning protocol, and with full tuning the kernel's margin over the network at the full block is
$1.8$ rather than $1.6$.

Full-data tuning barely changes the input-scaled kernel: $0.096\%$ against $0.095\%$, better on
four of ten splits, ratio $1.005$. Its coordinate search already arrives at the scale the full
block selects.

\paragraph{Gradient-based metric.} The coordinate search costs a sweep over the six inputs at
five multipliers, twice, per component and split. The alternative fixes the metric before any
kernel is fitted: the multiplier of coordinate $j$ is the root-mean-square gradient of the fitted
cubic surrogate along $j$ over the training rows, raised to a power $p$, normalized to geometric
mean one and clipped to $[1/16, 4]$, the range of the multipliers the search selected, so that
weakly varying coordinates receive longer length scales. It costs one cubic fit and six finite differences.

At $p=1$ the radiance error is $0.159\%$, below the isotropic kernel on every split, by a factor
$0.669$ on average, and above the searched metric on every split, by a factor $1.680$. The
dependence on $p$ is not monotone: $p=1/2$ gives $0.125\%$ and $p=1/4$ gives $0.157\%$. Choosing
$p$ like every other hyperparameter here, by the lowest mean validation error over the four
components, selects $p=1/2$ on all ten splits. At that exponent the gradient metric closes $79\%$
of the gap between the isotropic and searched kernels on radiance and $82\%$ on the median
reflectance error, for one cubic fit and a choice among three numbers.

The components prefer different exponents: $Y_2$ and $Y_3$ are best at $1/2$, the albedo at $1$
($2.17\%$ against $2.27\%$), and the path radiance at $1/4$. On path radiance the isotropic kernel
already reads $0.084\%$ and the searched metric $0.085\%$, while the proportional metric reads
$0.34\%$, four times the isotropic error. The gradient of the surrogate along a coordinate
measures how fast the map moves with it, not how much that coordinate helps predict, and on this
component the two disagree. Softening the exponent to $1/4$ brings $Y_1$ back to $0.085\%$, at the
cost of $Y_2$ and $Y_3$. Coordinate search can adjust each component separately, whereas the
gradient construction uses one exponent for all four.


\section{Comparison with OCO-2}\label{sec:oco2}

The companion paper \cite{nmkc} applies the same construction to the OCO-2 radiative-transfer
emulation problem of Lamminp\"a\"a et al.~\cite{lamminpaa}: per spectral band (O$_2$, weak
CO$_2$, strong CO$_2$), a reduced atmospheric state of twenty or twenty-four coordinates maps to
forty reduced radiance coefficients, with errors on the standardized coefficients and on the
reconstructed radiance, under the ten-split protocol of Section~\ref{sec:data}. It gives a
comparison at a higher input dimension.

There a kernel on the state trails the network by a factor of nine, four and five on the reduced
coefficients: the isotropic Mat\'ern reaches $40.6\%$, $58.6\%$ and $43.4\%$ on the three bands
against $4.4\%$, $16.4\%$ and $8.2\%$ for the network. Four ways of learning the kernel's metric
on the state were run on the same ten splits on all three bands: a kernel-flow metric
\cite{owhadiyoo}, a Mahalanobis metric under the same flow, the gradient outer product of the
recursive feature machine \cite{rfm}, and per-input scales fitted by empirical Bayes. The last is
the best of the four on every band, taking the isotropic kernel's $40.6\%$, $58.6\%$ and $44.3\%$
on the coefficients to $18.7\%$, $25.3\%$ and $16.9\%$, and on O$_2$ its radiance from $0.233\%$ to
$0.150\%$; the two flow metrics agree within half a point of each other at $26.9\%$, $29.8\%$ and
$23.9\%$, the recursive feature machine reads $28.2\%$, $47.9\%$ and $31.3\%$, and an additive
kernel $38.2\%$, $46.9\%$ and $47.7\%$. (The isotropic kernel on strong CO$_2$ reads $44.3\%$ here
and $43.4\%$ in the baseline experiment. The metric experiment also chooses the Mat\'ern smoothness
on validation and takes $\nu=3/2$ on this band at every split; the baseline fixes $\nu=5/2$, and its
choice of scale and nugget changes from split to split, six splits landing near $41\%$, three near
$48\%$ and one at $44\%$, hence its spread of $\pm3.4$.)
The feature kernel remains more accurate, by a margin that depends on the band: the exact Mat\'ern
regression on the network's frozen features reads $4.1\%$, $16.3\%$ and $8.1\%$ on the same splits,
so the best state metric is $4.6$ times that on O$_2$ but only $1.6$ and $2.1$ times it on the two
CO$_2$ bands. The empirical-Bayes metric is bimodal across splits on O$_2$, eight of ten landing near
$17.2\%$ and two at $24.7\%$, so its spread of $\pm3.2$ is thirty times the isotropic kernel's
$\pm0.1$. The exact Mat\'ern regression on the network's frozen features beats a ridge readout
refitted on the same features on every split, every band and every training loss, and a
per-coefficient choice among heads trained under different losses has lower mean error than the
release's own emulator under both reported metrics: it wins the reduced-coefficient comparison on
all ten splits on each band, and the radiance comparison on ten, nine and ten splits on O$_2$, weak
CO$_2$ and strong CO$_2$.

\section{Ten further emulation configurations}\label{sec:corpora}

Table~\ref{tab:corpora} extends the comparison to ten further configurations, each with
validation-based selection; the comparison concerns the families under a common protocol, not the
best method available for each problem.

The corpora are the advection equation with discontinuous initial data from the operator suite
of \cite{dehoop}, which is also the source of the structural-mechanics problem of \cite{nmkc};
Burgers' equation at three viscosities from PDEBench \cite{pdebench}; three one-step maps from
The Well \cite{thewell}: a two-dimensional turbulent radiative layer, an active-matter system and
a Helmholtz staircase; the ClimSim atmospheric physics emulation task \cite{climsim}; and the
paired correction-coefficient corpus of the physics-guided radiative-transfer work of \cite{pkan},
a third radiative-transfer emulation problem and the largest here, with about 488{,}000 training
rows. That corpus
is run in both of its configurations: from the atmospheric state alone, and with the low-fidelity
6S coefficients as additional inputs, which is a different and easier problem.

Advection is the 200-point input/output pair of the operator suite at full resolution. Burgers is
PDEBench's \texttt{1D\_Burgers\_Sols\_Nu}$\nu$ release, the map from the initial condition to the
solution at $t = 1$ on a grid subsampled by four to 256 cells, with the input used raw and the output
reduced to 128 principal coefficients. The three Well datasets are one-step maps of the full
multi-channel state with 256 principal coefficients in and out: the turbulent radiative layer on
its $128 \times 384$ grid (density, pressure, velocity), active matter on $256 \times 256$
(concentration, velocity and the two tensor fields), and the Helmholtz staircase on its two pressure
channels. ClimSim is the low-resolution release, 124 inputs to 128 tendencies, with 100{,}000
training rows drawn from the 10.1 million available and validation and test taken from the release's
own validation and scoring files. The correction-coefficient corpus is the paired 6S/libRadtran
Sentinel-2 set \cite{vermote,mayer}: wavelength, solar zenith, view zenith, relative azimuth,
aerosol optical depth, water vapor, ozone and elevation map to three libRadtran coefficients, split
by atmospheric state so that all bands of a state stay together; the multi-fidelity configuration
adds the three 6S coefficients, for input dimensions eight and eleven. For these two configurations
the error is
\begin{equation}\label{eq:coefficient-floor}
 E_{0.05}=\frac{1}{N}\sum_{i=1}^{N}
 \frac{\|Y_i-\widehat Y_i\|_2}{\max\{\|Y_i\|_2,0.05\}},
\end{equation}
and the same floor enters validation selection and the kernel-flow objective, which limits the
weight of coefficient vectors of small norm.

The families are those of Section~\ref{sec:methods} with two changes forced by scale. The
reference is a ridge regression linear in the inputs, since a cubic expansion of a 256-dimensional
input is not tractable. Kernel solves are exact when the training block has at most 20{,}000 rows
and use a Nystr\"om approximation \cite{nystrom} on 6{,}000 uniformly drawn landmarks above that,
which applies to ClimSim, the Helmholtz staircase at 20{,}384 rows and the two
correction-coefficient configurations. The convex stack and the per-coordinate selection are formed
from the kernels, the networks, the residual correction and the feature kernel; the ridge is a
reference row, not one of their members. Advection, Burgers, ClimSim and the two
correction-coefficient configurations have three to five random splits. The Well's releases carry
their own training, validation and test files, which are used unchanged, so the three
turbulent-radiative-layer runs share one split and differ only in network initialization (their
spread, of order $0.01$ points, measures initialization), and the active-matter and
Helmholtz-staircase rows are single runs.

\begin{table}[t]
\centering
% Generated by code/format_publication_tables.py; numerical cells preserved.
\textbf{A. Configurations}\par\medskip
{\fontsize{9}{11}\selectfont
\setlength{\tabcolsep}{4pt}
\begin{tabular}{lrrrr}
\toprule
corpus & $n$ & $d$ & $q$ & runs \\
\midrule
Advection & 19,000 & 200 & 200 & 5 \\
Burgers $\nu=0.1$ & 7,200 & 256 & 128 & 3 \\
Burgers $\nu=0.01$ & 7,200 & 256 & 128 & 3 \\
Burgers $\nu=0.001$ & 7,200 & 256 & 128 & 3 \\
Radiative layer & 7,200 & 256 & 256 & 3 \\
Active matter & 14,000 & 256 & 256 & 1 \\
Helmholtz staircase & 20,384 & 256 & 256 & 1 \\
ClimSim & 100,000 & 124 & 128 & 5 \\
RT coefficients, state & 487,773--487,823 & 8 & 3 & 5 \\
RT coefficients, +6S & 487,773--487,823 & 11 & 3 & 5 \\
\bottomrule
\end{tabular}}
\par\medskip\textbf{B. Reference, input kernels and network}\par\medskip
{\fontsize{9}{11}\selectfont
\setlength{\tabcolsep}{3pt}
\begin{tabular}{lrrrr}
\toprule
corpus & linear ridge & KRR & \thead{learned-metric\\KRR} & network \\
\midrule
Advection & 11.29$\pm$0.04 & 11.36$\pm$0.04 & 11.36$\pm$0.04 & 11.69$\pm$0.07 \\
Burgers $\nu=0.1$ & 13.89$\pm$0.53 & 2.81$\pm$0.28 & 2.90$\pm$0.33 & 20.34$\pm$0.61 \\
Burgers $\nu=0.01$ & 33.68$\pm$0.28 & 23.40$\pm$0.54 & 23.58$\pm$0.50 & 36.65$\pm$1.95 \\
Burgers $\nu=0.001$ & 38.38$\pm$0.12 & 31.66$\pm$0.16 & 33.23$\pm$1.73 & 39.98$\pm$0.65 \\
Radiative layer & 31.73 & 30.31 & 30.26$\pm$0.01 & 32.17$\pm$0.02 \\
Active matter & 37.24 & 35.81 & 35.18 & 37.31 \\
Helmholtz staircase & 0.00 & 0.26 & 0.07 & 0.40 \\
ClimSim & 32.24$\pm$0.06 & 27.20$\pm$0.05 & 26.31$\pm$0.04 & 25.10$\pm$0.18 \\
RT coefficients, state & 71.54$\pm$0.59 & 62.93$\pm$0.51 & 3.97$\pm$0.42 & 3.75$\pm$0.25 \\
RT coefficients, +6S & 33.83$\pm$0.26 & 15.76$\pm$1.98 & 4.60$\pm$0.31 & 2.79$\pm$0.14 \\
\bottomrule
\end{tabular}}
\par\medskip\textbf{C. Ensembles, corrections and combinations}\par\medskip
{\fontsize{9}{11}\selectfont
\setlength{\tabcolsep}{3pt}
\begin{tabular}{lrrrr}
\toprule
corpus & ensemble & \thead{network +\\residual KRR} & feature kernel & stack \\
\midrule
Advection & 11.50$\pm$0.05 & 11.50$\pm$0.05 & 11.64$\pm$0.07 & 11.36$\pm$0.04 \\
Burgers $\nu=0.1$ & 19.28$\pm$0.09 & 2.78$\pm$0.29 & 4.92$\pm$0.30 & 2.78$\pm$0.29 \\
Burgers $\nu=0.01$ & 35.31$\pm$0.70 & 23.12$\pm$0.56 & 29.20$\pm$0.59 & 23.01$\pm$0.54 \\
Burgers $\nu=0.001$ & 38.89$\pm$0.11 & 31.82$\pm$0.11 & 35.59$\pm$0.48 & 31.34$\pm$0.14 \\
Radiative layer & 31.58$\pm$0.02 & 30.40$\pm$0.01 & 31.08$\pm$0.01 & 30.12 \\
Active matter & 36.93 & 35.77 & 36.22 & 35.05 \\
Helmholtz staircase & 0.24 & 0.23 & 0.45 & 0.07 \\
ClimSim & 24.43$\pm$0.17 & 24.33$\pm$0.13 & 24.71$\pm$0.25 & 24.25$\pm$0.17 \\
RT coefficients, state & 3.24$\pm$0.11 & 3.13$\pm$0.05 & 3.52$\pm$0.07 & 3.09$\pm$0.05 \\
RT coefficients, +6S & 2.40$\pm$0.11 & 2.34$\pm$0.09 & 2.80$\pm$0.09 & 2.32$\pm$0.08 \\
\bottomrule
\end{tabular}}

\caption{Ten further configurations: dimensions and run counts (A), reference models (B),
and combinations (C). Errors are mean relative spectral norms in percent, with the
0.05 denominator floor of \eqref{eq:coefficient-floor} for the two RT-coefficient rows.
Spreads are sample standard deviations; the radiative-layer repetitions share one split
and vary initialization, and two rows are single runs. The network budget is 200 epochs
at width 512, except ClimSim (60 at 512) and RT coefficients (120 at 384). The linear
ridge is excluded from the stack; its Helmholtz error is $0.004\%$, rounded in the table.
Kernel solves above 20{,}000 training rows use 6{,}000 Nystr\"om landmarks.}
\label{tab:corpora}
\end{table}

The relative errors vary widely with the problem. Where the map is smooth and the sample small,
Burgers at $\nu = 0.1$, the exact kernel is at $2.81\%$ against the network's $20.34\%$, a factor of
seven, of the same kind as the five- to six-fold gap between the per-input kernel and the network on
the EMIT table's smooth components; as the viscosity falls and the solutions roughen the gap closes
to $23.40\%$ against $36.65\%$ and then to $31.66\%$ against $39.98\%$. Where the sample is large and
the map rough, on ClimSim and on the correction-coefficient corpus at nearly half a million rows, the
network leads instead, $25.10\%$ against $27.20\%$ and $3.75\%$ against $62.93\%$, which is the OCO-2
regime. On advection, where the operator is linear, the linear ridge is the lowest row and every
other family lands within $0.4$ points of it.

Learned metrics help most on the correction-coefficient corpus. From the state alone the isotropic
kernel is at $62.93 \pm 0.51$ and a kernel-flow metric \cite{owhadiyoo} takes it to $3.97 \pm 0.42$,
a factor of sixteen, which brings a kernel on the state within six percent of the network on a
corpus where the plain kernel was seventeen times worse. With the low-fidelity coefficients as
inputs the same pair reads $15.76$ and $4.60$, a factor of $3.4$: the easier the problem is made for
the plain kernel, the less the learned metric adds.

The convex stack is the lowest row on eight of the ten configurations and the second lowest on the
other two; the residual-corrected network is among the lowest two on five. Both exceptions are the
linear problems, and on both the row that beats the stack is the ridge, which the stack does not
contain. On advection the ridge is lower by $0.07$ points while the stack, at $11.36\%$, matches its
best member, the isotropic kernel, to two decimals; that kernel and the learned-metric kernel, which
differ by less than $0.01$ points there, carry between $0.85$ and $0.91$ of its weight at every
split. On the Helmholtz
staircase the ridge reaches $0.004\%$, sixteen times below the next row and about a hundred times
below the weakest member of the pool, and the stack reaches $0.066\%$, slightly below its best member,
the learned-metric kernel at $0.067\%$, on which it puts weight $0.94$. In both cases the stack does
what it is asked within its pool and cannot recover a family the pool excludes. A combination is a
good default when no member of its pool dominates it, and it can be read against the whole
comparison only when the pool contains the families being compared.

These rows use one training budget per corpus and one architecture family, and several of the
corpora have published results from methods tuned for them. With three to five runs, and one on two
rows, the spreads support the large orderings above and not fine distinctions between adjacent
families. Every kernel column also inherits the tuning protocol of Section~\ref{sec:retune}, with
the length scale chosen on a $6{,}000$-row subsample and the winner refit on every row. On the EMIT
table, choosing it with every row in the solve lowers the isotropic kernel's radiance error by a
factor $0.882$ on average, so the kernel columns of Table~\ref{tab:corpora} describe the stated
training and selection procedures rather than the best each kernel can reach.


\clearpage
\section{Extended diagnostic tables}\label{app:diagnostics}

This section collects the perturbation, ridge-path, synthetic and external-benchmark tables referred to in
Sections~\ref{sec:results-representation}--\ref{sec:results-benchmarks}, and the record of the runs behind
Section~\ref{sec:tq-policy} of the article.

\begin{table}[h]
\centering
% generated by code/make_tq_tables.py; seeds w512 [101, 102, 103, 106, 107, 108] w2000 [101, 102, 103, 106, 107, 108]
\begin{tabular}{lllrrrrr}
\toprule
width & seed & machine & raw & admissible & matched & validation & test \\
\midrule
3$\times$512 & 101 & DGX & 18,884 & 15,287 & 15,287 & 2,098 & 2,331 \\
 & 102 & DGX & 18,884 & 15,239 & 15,239 & 2,098 & 2,331 \\
 & 103 & DGX & 18,884 & 15,284 & 15,284 & 2,098 & 2,331 \\
 & 106 & Kaggle & 18,884 & 15,290 & 15,290 & 2,098 & 2,331 \\
 & 107 & Kaggle & 18,884 & 15,325 & 15,325 & 2,098 & 2,331 \\
 & 108 & Kaggle & 18,884 & 15,250 & 15,250 & 2,098 & 2,331 \\
3$\times$2000 & 101 & DGX & 18,884 & 15,287 & 15,287 & 2,098 & 2,331 \\
 & 102 & DGX & 18,884 & 15,239 & 15,239 & 2,098 & 2,331 \\
 & 103 & DGX & 18,884 & 15,284 & 15,284 & 2,098 & 2,331 \\
 & 106 & Kaggle & 18,884 & 15,290 & 15,290 & 2,098 & 2,331 \\
 & 107 & Kaggle & 18,884 & 15,325 & 15,325 & 2,098 & 2,331 \\
 & 108 & Kaggle & 18,884 & 15,250 & 15,250 & 2,098 & 2,331 \\
\bottomrule
\end{tabular}

\caption{The runs of the training-target experiment: for each split and width, the machine it ran on
and the number of training states in each arm. Validation and test blocks are the same in the three
arms. The admissible and matched arms have the same size by construction.}
\label{tab:tq-runs}
\end{table}

% Generated by code/format_publication_tables.py; numerical cells preserved.
\clearpage
\begin{landscape}
{\fontsize{9}{11}\selectfont
\setlength{\tabcolsep}{4pt}
\setlength{\LTcapwidth}{0.95\linewidth}
\begin{longtable}{llrrrrrrr}
\caption{Same-data perturbation diagnostics for Proposition~\ref{prop:perturb}. Movement is a Euclidean prediction change in retained coordinates. The four ratio columns are \emph{medians of entrywise bound/movement ratios}, not ratios of medians. The fourteen state-input rows are retained historical values whose JSON is absent. The feature rows are regenerated from four component records at seed 101 and four at seed 104, each with 4,000 fitting and 1,000 query rows. The signed minimum eigenvalue is rounded: a displayed $-1.000$ does not assert exact equality to $-1$. A dash means the corresponding bound is not applicable.}\label{tab:v2-perturbation}\\
\toprule
source & perturbation & $\lambda_{\min}(B)$ & $\|B\|$ & \thead{median\\movement} & \thead{residual-action\\ratio} & \thead{signed-spectrum\\ratio} & \thead{$\|B\|<1$\\ratio} & \thead{ridge\\ratio} \\
\midrule
\endfirsthead
\multicolumn{9}{l}{\tablename\ \thetable\ (continued)}\\
\toprule
source & perturbation & $\lambda_{\min}(B)$ & $\|B\|$ & \thead{median\\movement} & \thead{residual-action\\ratio} & \thead{signed-spectrum\\ratio} & \thead{$\|B\|<1$\\ratio} & \thead{ridge\\ratio} \\
\midrule
\endhead
\midrule
\multicolumn{9}{r}{Continued on next page}\\
\endfoot
\bottomrule
\endlastfoot
state (archive) & Nystr\"om 50 & -0.941 & 0.941 & 2.01 & 57.3 & 4.05e+03 & 4.05e+03 & -- \\
state (archive) & Nystr\"om 100 & -0.872 & 0.872 & 1.05 & 82.1 & 3.05e+03 & 3.05e+03 & -- \\
state (archive) & Nystr\"om 200 & -0.642 & 0.642 & 0.498 & 136 & 1.65e+03 & 1.65e+03 & -- \\
state (archive) & Nystr\"om 400 & -0.411 & 0.411 & 0.2 & 178 & 1.54e+03 & 1.54e+03 & -- \\
state (archive) & Nystr\"om 800 & -0.193 & 0.193 & 0.0702 & 264 & 1.47e+03 & 1.47e+03 & -- \\
state (archive) & Nystr\"om 1600 & -0.078 & 0.078 & 0.0103 & 408 & 3.51e+03 & 3.51e+03 & -- \\
state (archive) & ridge $\times$0.5 & -0.500 & 0.500 & 0.304 & 517 & 1.34e+03 & 1.34e+03 & -- \\
state (archive) & ridge $\times$2.0 & +0.000 & 0.999 & 0.372 & 835 & 550 & 1.32e+06 & 550 \\
state (archive) & ridge $\times$10.0 & +0.002 & 8.993 & 1.36 & 1.98e+03 & 271 & -- & 271 \\
state (archive) & ridge $\times$100.0 & +0.018 & 98.918 & 3.67 & 7.38e+03 & 111 & -- & 111 \\
state (archive) & PCA inputs 2 & -0.998 & 3.600 & 10.5 & 14.3 & 4.81e+04 & -- & -- \\
state (archive) & PCA inputs 3 & -0.998 & 2.523 & 9.8 & 14.2 & 4.49e+04 & -- & -- \\
state (archive) & PCA inputs 4 & -0.998 & 1.581 & 5.24 & 25.3 & 6.52e+04 & -- & -- \\
state (archive) & PCA inputs 5 & -0.997 & 0.997 & 1.9 & 64.3 & 1.39e+05 & 1.39e+05 & -- \\
s101, $Y_1$ & PCA features 1000 & -0.192 & 0.192 & 0.00383 & 4.01e+03 & 3.8e+04 & 3.8e+04 & -- \\
s101, $Y_1$ & PCA features 500 & -0.503 & 0.503 & 0.0107 & 4.41e+03 & 5.82e+04 & 5.82e+04 & -- \\
s101, $Y_1$ & PCA features 250 & -0.764 & 0.764 & 0.0185 & 4.3e+03 & 1.1e+05 & 1.1e+05 & -- \\
s101, $Y_1$ & PCA features 125 & -0.917 & 0.917 & 0.0279 & 3.88e+03 & 2.6e+05 & 2.6e+05 & -- \\
s101, $Y_1$ & PCA features 64 & -0.965 & 0.965 & 0.0407 & 3.55e+03 & 4.78e+05 & 4.78e+05 & -- \\
s101, $Y_1$ & PCA features 32 & -0.991 & 0.991 & 0.0595 & 3.13e+03 & 1.59e+06 & 1.59e+06 & -- \\
s101, $Y_1$ & PCA features 16 & -0.999 & 0.999 & 0.0923 & 2.48e+03 & 1.3e+07 & 1.3e+07 & -- \\
s101, $Y_1$ & Nystr\"om 100 & -1.000 & 1.000 & 0.465 & 637 & 8.85e+07 & 8.85e+07 & -- \\
s101, $Y_1$ & Nystr\"om 400 & -1.000 & 1.000 & 0.223 & 1.27e+03 & 7.11e+07 & 7.11e+07 & -- \\
s101, $Y_1$ & Nystr\"om 1600 & -1.000 & 1.000 & 0.0746 & 2.58e+03 & 2.74e+07 & 2.74e+07 & -- \\
s101, $Y_1$ & Nystr\"om 3200 & -0.997 & 0.997 & 0.0196 & 5.22e+03 & 9.75e+06 & 9.75e+06 & -- \\
s101, $Y_1$ & ridge $\times$0.5 & -0.490 & 0.490 & 0.0102 & 9.61e+03 & 5.74e+04 & 5.74e+04 & -- \\
s101, $Y_1$ & ridge $\times$2.0 & +0.000 & 0.979 & 0.0115 & 1.7e+04 & 2.62e+04 & 2.5e+06 & 2.62e+04 \\
s101, $Y_1$ & ridge $\times$10.0 & +0.000 & 8.813 & 0.0407 & 4.32e+04 & 1.34e+04 & -- & 1.34e+04 \\
s101, $Y_1$ & ridge $\times$100.0 & +0.000 & 96.942 & 0.0889 & 2.17e+05 & 6.78e+03 & -- & 6.78e+03 \\
s101, $Y_2$ & PCA features 1000 & -0.122 & 0.122 & 0.00143 & 6.41e+03 & 4.95e+04 & 4.95e+04 & -- \\
s101, $Y_2$ & PCA features 500 & -0.418 & 0.418 & 0.00425 & 6.8e+03 & 8.59e+04 & 8.59e+04 & -- \\
s101, $Y_2$ & PCA features 250 & -0.726 & 0.726 & 0.00802 & 6.45e+03 & 1.7e+05 & 1.7e+05 & -- \\
s101, $Y_2$ & PCA features 125 & -0.868 & 0.868 & 0.0117 & 6.33e+03 & 2.99e+05 & 2.99e+05 & -- \\
s101, $Y_2$ & PCA features 64 & -0.977 & 0.977 & 0.0184 & 5.76e+03 & 1.3e+06 & 1.3e+06 & -- \\
s101, $Y_2$ & PCA features 32 & -0.989 & 0.989 & 0.0267 & 5.78e+03 & 2.3e+06 & 2.3e+06 & -- \\
s101, $Y_2$ & PCA features 16 & -0.999 & 0.999 & 0.0419 & 4.83e+03 & 1.49e+07 & 1.49e+07 & -- \\
s101, $Y_2$ & Nystr\"om 100 & -1.000 & 1.000 & 0.345 & 731 & 1.01e+08 & 1.01e+08 & -- \\
s101, $Y_2$ & Nystr\"om 400 & -1.000 & 1.000 & 0.122 & 1.97e+03 & 8.48e+07 & 8.48e+07 & -- \\
s101, $Y_2$ & Nystr\"om 1600 & -1.000 & 1.000 & 0.0823 & 2.38e+03 & 5.34e+07 & 5.34e+07 & -- \\
s101, $Y_2$ & Nystr\"om 3200 & -0.998 & 0.998 & 0.0161 & 7.05e+03 & 1.86e+07 & 1.86e+07 & -- \\
s101, $Y_2$ & ridge $\times$0.5 & -0.491 & 0.491 & 0.00632 & 1.71e+04 & 7.71e+04 & 7.71e+04 & -- \\
s101, $Y_2$ & ridge $\times$2.0 & +0.000 & 0.982 & 0.00687 & 3.14e+04 & 3.64e+04 & 3.96e+06 & 3.64e+04 \\
s101, $Y_2$ & ridge $\times$10.0 & +0.000 & 8.836 & 0.0249 & 7.8e+04 & 1.82e+04 & -- & 1.82e+04 \\
s101, $Y_2$ & ridge $\times$100.0 & +0.000 & 97.199 & 0.0545 & 3.92e+05 & 9.17e+03 & -- & 9.17e+03 \\
s101, $Y_3$ & PCA features 1000 & -0.170 & 0.170 & 0.00292 & 6.17e+03 & 4.56e+04 & 4.56e+04 & -- \\
s101, $Y_3$ & PCA features 500 & -0.475 & 0.475 & 0.00766 & 7.02e+03 & 7.73e+04 & 7.73e+04 & -- \\
s101, $Y_3$ & PCA features 250 & -0.738 & 0.738 & 0.0136 & 6.44e+03 & 1.38e+05 & 1.38e+05 & -- \\
s101, $Y_3$ & PCA features 125 & -0.898 & 0.898 & 0.0211 & 5.61e+03 & 2.87e+05 & 2.87e+05 & -- \\
s101, $Y_3$ & PCA features 64 & -0.964 & 0.964 & 0.028 & 5.45e+03 & 7.06e+05 & 7.06e+05 & -- \\
s101, $Y_3$ & PCA features 32 & -0.990 & 0.990 & 0.0406 & 4.9e+03 & 2.09e+06 & 2.09e+06 & -- \\
s101, $Y_3$ & PCA features 16 & -0.999 & 0.999 & 0.0725 & 3.51e+03 & 1.63e+07 & 1.63e+07 & -- \\
s101, $Y_3$ & Nystr\"om 100 & -1.000 & 1.000 & 0.35 & 876 & 4.97e+07 & 4.97e+07 & -- \\
s101, $Y_3$ & Nystr\"om 400 & -1.000 & 1.000 & 0.191 & 1.55e+03 & 4.53e+07 & 4.53e+07 & -- \\
s101, $Y_3$ & Nystr\"om 1600 & -0.999 & 0.999 & 0.0709 & 3.33e+03 & 1.45e+07 & 1.45e+07 & -- \\
s101, $Y_3$ & Nystr\"om 3200 & -0.997 & 0.997 & 0.0331 & 4.54e+03 & 6.39e+06 & 6.39e+06 & -- \\
s101, $Y_3$ & ridge $\times$0.5 & -0.492 & 0.492 & 0.00865 & 1.65e+04 & 7.26e+04 & 7.26e+04 & -- \\
s101, $Y_3$ & ridge $\times$2.0 & +0.000 & 0.985 & 0.00949 & 3e+04 & 3.38e+04 & 4.43e+06 & 3.38e+04 \\
s101, $Y_3$ & ridge $\times$10.0 & +0.000 & 8.864 & 0.0357 & 7.18e+04 & 1.63e+04 & -- & 1.63e+04 \\
s101, $Y_3$ & ridge $\times$100.0 & +0.000 & 97.500 & 0.0821 & 3.44e+05 & 7.81e+03 & -- & 7.81e+03 \\
s101, $Y_4$ & PCA features 1000 & -0.006 & 0.006 & 0.00138 & 538 & 4.64e+03 & 4.64e+03 & -- \\
s101, $Y_4$ & PCA features 500 & -0.032 & 0.032 & 0.00463 & 712 & 7.83e+03 & 7.83e+03 & -- \\
s101, $Y_4$ & PCA features 250 & -0.114 & 0.114 & 0.0104 & 826 & 1.35e+04 & 1.35e+04 & -- \\
s101, $Y_4$ & PCA features 125 & -0.269 & 0.269 & 0.0191 & 897 & 2.12e+04 & 2.12e+04 & -- \\
s101, $Y_4$ & PCA features 64 & -0.594 & 0.594 & 0.0286 & 963 & 5.8e+04 & 5.8e+04 & -- \\
s101, $Y_4$ & PCA features 32 & -0.741 & 0.741 & 0.0436 & 983 & 8.06e+04 & 8.06e+04 & -- \\
s101, $Y_4$ & PCA features 16 & -0.902 & 0.902 & 0.0769 & 987 & 1.65e+05 & 1.65e+05 & -- \\
s101, $Y_4$ & Nystr\"om 100 & -0.998 & 0.998 & 0.603 & 304 & 1.2e+06 & 1.2e+06 & -- \\
s101, $Y_4$ & Nystr\"om 400 & -0.985 & 0.985 & 0.35 & 500 & 2.11e+05 & 2.11e+05 & -- \\
s101, $Y_4$ & Nystr\"om 1600 & -0.860 & 0.860 & 0.135 & 916 & 4.86e+04 & 4.86e+04 & -- \\
s101, $Y_4$ & Nystr\"om 3200 & -0.762 & 0.762 & 0.0397 & 1.71e+03 & 8.46e+04 & 8.46e+04 & -- \\
s101, $Y_4$ & ridge $\times$0.5 & -0.494 & 0.494 & 0.063 & 3.15e+03 & 1.61e+04 & 1.61e+04 & -- \\
s101, $Y_4$ & ridge $\times$2.0 & +0.000 & 0.989 & 0.0599 & 6.6e+03 & 8.63e+03 & 1.56e+06 & 8.63e+03 \\
s101, $Y_4$ & ridge $\times$10.0 & +0.000 & 8.901 & 0.191 & 1.86e+04 & 4.89e+03 & -- & 4.89e+03 \\
s101, $Y_4$ & ridge $\times$100.0 & +0.000 & 97.908 & 0.36 & 1.07e+05 & 2.86e+03 & -- & 2.86e+03 \\
s104, $Y_1$ & PCA features 1000 & -0.184 & 0.184 & 0.00397 & 4.68e+03 & 3.85e+04 & 3.85e+04 & -- \\
s104, $Y_1$ & PCA features 500 & -0.471 & 0.471 & 0.0104 & 5.01e+03 & 5.82e+04 & 5.82e+04 & -- \\
s104, $Y_1$ & PCA features 250 & -0.724 & 0.724 & 0.0175 & 4.8e+03 & 1.04e+05 & 1.04e+05 & -- \\
s104, $Y_1$ & PCA features 125 & -0.901 & 0.901 & 0.0267 & 4.25e+03 & 2.45e+05 & 2.45e+05 & -- \\
s104, $Y_1$ & PCA features 64 & -0.959 & 0.959 & 0.0373 & 3.87e+03 & 4.76e+05 & 4.76e+05 & -- \\
s104, $Y_1$ & PCA features 32 & -0.985 & 0.985 & 0.055 & 3.43e+03 & 1.07e+06 & 1.07e+06 & -- \\
s104, $Y_1$ & PCA features 16 & -0.997 & 0.997 & 0.0893 & 2.61e+03 & 3.51e+06 & 3.51e+06 & -- \\
s104, $Y_1$ & concatenation & -0.749 & 76.009 & 0.0427 & 2.3e+05 & 1.99e+05 & -- & -- \\
s104, $Y_1$ & member 1 & -0.973 & 175.948 & 0.049 & 2.34e+05 & 2.24e+06 & -- & -- \\
s104, $Y_1$ & member 2 & -0.967 & 126.717 & 0.0541 & 3.81e+05 & 2.02e+06 & -- & -- \\
s104, $Y_1$ & Nystr\"om 100 & -1.000 & 1.000 & 0.861 & 379 & 9.61e+07 & 9.61e+07 & -- \\
s104, $Y_1$ & Nystr\"om 400 & -1.000 & 1.000 & 0.247 & 1.25e+03 & 6.69e+07 & 6.69e+07 & -- \\
s104, $Y_1$ & Nystr\"om 1600 & -1.000 & 1.000 & 0.112 & 2.01e+03 & 4.7e+07 & 4.7e+07 & -- \\
s104, $Y_1$ & Nystr\"om 3200 & -1.000 & 1.000 & 0.0249 & 4.43e+03 & 5.86e+07 & 5.86e+07 & -- \\
s104, $Y_1$ & ridge $\times$0.5 & -0.493 & 0.493 & 0.011 & 1.06e+04 & 5.96e+04 & 5.96e+04 & -- \\
s104, $Y_1$ & ridge $\times$2.0 & +0.000 & 0.986 & 0.0124 & 1.89e+04 & 2.7e+04 & 3.93e+06 & 2.7e+04 \\
s104, $Y_1$ & ridge $\times$10.0 & +0.000 & 8.877 & 0.0445 & 4.73e+04 & 1.36e+04 & -- & 1.36e+04 \\
s104, $Y_1$ & ridge $\times$100.0 & +0.000 & 97.646 & 0.102 & 2.27e+05 & 6.54e+03 & -- & 6.54e+03 \\
s104, $Y_2$ & PCA features 1000 & -0.108 & 0.108 & 0.00163 & 7.72e+03 & 5.02e+04 & 5.02e+04 & -- \\
s104, $Y_2$ & PCA features 500 & -0.370 & 0.370 & 0.00553 & 8.84e+03 & 7.23e+04 & 7.23e+04 & -- \\
s104, $Y_2$ & PCA features 250 & -0.668 & 0.668 & 0.0107 & 8.77e+03 & 1.29e+05 & 1.29e+05 & -- \\
s104, $Y_2$ & PCA features 125 & -0.834 & 0.834 & 0.0166 & 7.71e+03 & 2.12e+05 & 2.12e+05 & -- \\
s104, $Y_2$ & PCA features 64 & -0.955 & 0.955 & 0.0234 & 6.59e+03 & 6.62e+05 & 6.62e+05 & -- \\
s104, $Y_2$ & PCA features 32 & -0.983 & 0.983 & 0.0326 & 6.12e+03 & 1.46e+06 & 1.46e+06 & -- \\
s104, $Y_2$ & PCA features 16 & -0.998 & 0.998 & 0.0503 & 5.04e+03 & 1.05e+07 & 1.05e+07 & -- \\
s104, $Y_2$ & concatenation & -0.734 & 96.255 & 0.0279 & 6.7e+05 & 3.68e+05 & -- & -- \\
s104, $Y_2$ & member 1 & -0.974 & 200.331 & 0.0328 & 8.4e+05 & 5.59e+06 & -- & -- \\
s104, $Y_2$ & member 2 & -0.968 & 266.674 & 0.035 & 9.16e+05 & 3.4e+06 & -- & -- \\
s104, $Y_2$ & Nystr\"om 100 & -1.000 & 1.000 & 0.713 & 451 & 1.85e+08 & 1.85e+08 & -- \\
s104, $Y_2$ & Nystr\"om 400 & -1.000 & 1.000 & 0.119 & 2.5e+03 & 4.07e+07 & 4.07e+07 & -- \\
s104, $Y_2$ & Nystr\"om 1600 & -0.999 & 0.999 & 0.0322 & 6.59e+03 & 1.8e+07 & 1.8e+07 & -- \\
s104, $Y_2$ & Nystr\"om 3200 & -0.992 & 0.992 & 0.0104 & 1.25e+04 & 8.1e+06 & 8.1e+06 & -- \\
s104, $Y_2$ & ridge $\times$0.5 & -0.498 & 0.498 & 0.00788 & 2.13e+04 & 8.49e+04 & 8.49e+04 & -- \\
s104, $Y_2$ & ridge $\times$2.0 & +0.000 & 0.995 & 0.00862 & 3.89e+04 & 3.91e+04 & 1.66e+07 & 3.91e+04 \\
s104, $Y_2$ & ridge $\times$10.0 & +0.000 & 8.958 & 0.0323 & 9.36e+04 & 1.88e+04 & -- & 1.88e+04 \\
s104, $Y_2$ & ridge $\times$100.0 & +0.000 & 98.534 & 0.0733 & 4.53e+05 & 9.13e+03 & -- & 9.13e+03 \\
s104, $Y_3$ & PCA features 1000 & -0.176 & 0.176 & 0.00305 & 6.26e+03 & 4.48e+04 & 4.48e+04 & -- \\
s104, $Y_3$ & PCA features 500 & -0.474 & 0.474 & 0.00801 & 7.06e+03 & 7.26e+04 & 7.26e+04 & -- \\
s104, $Y_3$ & PCA features 250 & -0.738 & 0.738 & 0.0145 & 6.18e+03 & 1.26e+05 & 1.26e+05 & -- \\
s104, $Y_3$ & PCA features 125 & -0.882 & 0.882 & 0.0189 & 6.1e+03 & 2.67e+05 & 2.67e+05 & -- \\
s104, $Y_3$ & PCA features 64 & -0.957 & 0.957 & 0.0259 & 5.75e+03 & 6.3e+05 & 6.3e+05 & -- \\
s104, $Y_3$ & PCA features 32 & -0.990 & 0.990 & 0.0371 & 5.12e+03 & 2.16e+06 & 2.16e+06 & -- \\
s104, $Y_3$ & PCA features 16 & -0.999 & 0.999 & 0.0648 & 3.65e+03 & 2.51e+07 & 2.51e+07 & -- \\
s104, $Y_3$ & concatenation & -0.749 & 12.229 & 0.0228 & 6.41e+04 & 2.52e+05 & -- & -- \\
s104, $Y_3$ & member 1 & -0.987 & 28.650 & 0.0374 & 1.01e+05 & 5.18e+06 & -- & -- \\
s104, $Y_3$ & member 2 & -0.978 & 32.067 & 0.0391 & 3.02e+04 & 2.59e+06 & -- & -- \\
s104, $Y_3$ & Nystr\"om 100 & -1.000 & 1.000 & 0.777 & 389 & 1.27e+08 & 1.27e+08 & -- \\
s104, $Y_3$ & Nystr\"om 400 & -1.000 & 1.000 & 0.203 & 1.42e+03 & 5.4e+07 & 5.4e+07 & -- \\
s104, $Y_3$ & Nystr\"om 1600 & -1.000 & 1.000 & 0.134 & 1.72e+03 & 4.38e+07 & 4.38e+07 & -- \\
s104, $Y_3$ & Nystr\"om 3200 & -0.999 & 0.999 & 0.029 & 4.4e+03 & 1.58e+07 & 1.58e+07 & -- \\
s104, $Y_3$ & ridge $\times$0.5 & -0.487 & 0.487 & 0.00833 & 1.62e+04 & 7.25e+04 & 7.25e+04 & -- \\
s104, $Y_3$ & ridge $\times$2.0 & +0.000 & 0.975 & 0.00917 & 2.93e+04 & 3.42e+04 & 2.65e+06 & 3.42e+04 \\
s104, $Y_3$ & ridge $\times$10.0 & +0.000 & 8.771 & 0.0329 & 7.36e+04 & 1.73e+04 & -- & 1.73e+04 \\
s104, $Y_3$ & ridge $\times$100.0 & +0.000 & 96.479 & 0.0765 & 3.48e+05 & 8.23e+03 & -- & 8.23e+03 \\
s104, $Y_4$ & PCA features 1000 & -0.003 & 0.003 & 0.000765 & 496 & 4.96e+03 & 4.96e+03 & -- \\
s104, $Y_4$ & PCA features 500 & -0.019 & 0.019 & 0.00303 & 682 & 8.81e+03 & 8.81e+03 & -- \\
s104, $Y_4$ & PCA features 250 & -0.088 & 0.088 & 0.00823 & 778 & 1.58e+04 & 1.58e+04 & -- \\
s104, $Y_4$ & PCA features 125 & -0.259 & 0.259 & 0.0153 & 851 & 3.11e+04 & 3.11e+04 & -- \\
s104, $Y_4$ & PCA features 64 & -0.406 & 0.406 & 0.0238 & 906 & 4e+04 & 4e+04 & -- \\
s104, $Y_4$ & PCA features 32 & -0.790 & 0.790 & 0.0427 & 859 & 1.27e+05 & 1.27e+05 & -- \\
s104, $Y_4$ & PCA features 16 & -0.883 & 0.883 & 0.0733 & 851 & 1.62e+05 & 1.62e+05 & -- \\
s104, $Y_4$ & concatenation & -0.658 & 58.594 & 0.0605 & 1.49e+04 & 1.97e+05 & -- & -- \\
s104, $Y_4$ & member 1 & -0.897 & 23.688 & 0.103 & 5.19e+03 & 3.18e+05 & -- & -- \\
s104, $Y_4$ & member 2 & -0.919 & 155.529 & 0.108 & 1.74e+04 & 8.07e+05 & -- & -- \\
s104, $Y_4$ & Nystr\"om 100 & -0.993 & 0.993 & 0.517 & 378 & 3.93e+05 & 3.93e+05 & -- \\
s104, $Y_4$ & Nystr\"om 400 & -0.961 & 0.961 & 0.241 & 772 & 1.39e+05 & 1.39e+05 & -- \\
s104, $Y_4$ & Nystr\"om 1600 & -0.850 & 0.850 & 0.0775 & 1.66e+03 & 9.85e+04 & 9.85e+04 & -- \\
s104, $Y_4$ & Nystr\"om 3200 & -0.707 & 0.707 & 0.0213 & 2.57e+03 & 1.5e+05 & 1.5e+05 & -- \\
s104, $Y_4$ & ridge $\times$0.5 & -0.491 & 0.491 & 0.0482 & 4.89e+03 & 2.63e+04 & 2.63e+04 & -- \\
s104, $Y_4$ & ridge $\times$2.0 & +0.000 & 0.983 & 0.0525 & 8.97e+03 & 1.24e+04 & 1.43e+06 & 1.24e+04 \\
s104, $Y_4$ & ridge $\times$10.0 & +0.000 & 8.846 & 0.179 & 2.36e+04 & 6.6e+03 & -- & 6.6e+03 \\
s104, $Y_4$ & ridge $\times$100.0 & +0.000 & 97.301 & 0.393 & 1.18e+05 & 3.32e+03 & -- & 3.32e+03 \\
\end{longtable}}
\end{landscape}

% Generated by code/format_publication_tables.py; numerical cells preserved.
\clearpage
\begin{landscape}
{\fontsize{9}{11}\selectfont
\setlength{\tabcolsep}{4pt}
\setlength{\LTcapwidth}{0.95\linewidth}
\begin{longtable}{llrrrrrrrrr}
\caption{Ridge-path diagnostics, seed 101, 4,000 fitting and 1,000 query rows. The bias and noise factor refer to Proposition~\ref{prop:sensitivity}(iii); the factor multiplies a hypothetical noise variance and is not an estimate of simulator noise. Corrected and kernel-only test columns are mean relative errors of a physical spectral component, in percent. The final column uses retained PCA coordinates on the held-out half and compares same-half fitting with split-half fitting; it is a different metric and sample. Each ridge is evaluated descriptively, not selected on these test values.}\label{tab:v2-ridge}\\
\toprule
component & kernel & $\gamma$ & $\mathrm{tr}S_\lambda$ & bias & \thead{noise\\factor} & \thead{median\\$\|a_\lambda\|^2$} & \thead{median\\$\widetilde P_\lambda$} & \thead{corrected\\{[\%]}} & \thead{kernel only\\{[\%]}} & \thead{same / split half\\PCA error [\%]} \\
\midrule
\endfirsthead
\multicolumn{11}{l}{\tablename\ \thetable\ (continued)}\\
\toprule
component & kernel & $\gamma$ & $\mathrm{tr}S_\lambda$ & bias & \thead{noise\\factor} & \thead{median\\$\|a_\lambda\|^2$} & \thead{median\\$\widetilde P_\lambda$} & \thead{corrected\\{[\%]}} & \thead{kernel only\\{[\%]}} & \thead{same / split half\\PCA error [\%]} \\
\midrule
\endhead
\midrule
\multicolumn{11}{r}{Continued on next page}\\
\endfoot
\bottomrule
\endlastfoot
Y1 & input & 1e-08 & 3841 & 0.000226 & 0.924 & 0.87 & 0.0324 & 0.2168 & 0.1754 & 0.079 / 2.465 \\
Y1 & input & 1e-06 & 1624 & 0.0712 & 0.26 & 0.226 & 0.0395 & 0.2376 & 0.2420 & 1.145 / 2.214 \\
Y1 & input & 0.0001 & 212 & 0.282 & 0.0269 & 0.0237 & 0.0997 & 0.4132 & 0.9768 & 2.221 / 2.395 \\
Y1 & input & 0.01 & 19 & 0.413 & 0.00194 & 0.00197 & 0.315 & 0.5095 & 5.7578 & 2.716 / 2.728 \\
Y1 & feature & 1e-08 & 2541 & 0.0012 & 0.51 & 0.481 & 0.00637 & 0.2027 & 0.2111 & 0.280 / 1.171 \\
Y1 & feature & 1e-06 & 689 & 0.0225 & 0.111 & 0.0668 & 0.0151 & 0.3130 & 0.3275 & 1.085 / 1.464 \\
Y1 & feature & 0.0001 & 103 & 0.156 & 0.0137 & 0.00825 & 0.0529 & 0.4546 & 0.6396 & 2.105 / 2.187 \\
Y1 & feature & 0.01 & 9 & 0.418 & 0.001 & 0.000845 & 0.185 & 0.5326 & 5.9929 & 2.783 / 2.786 \\
Y2 & input & 1e-08 & 3884 & 0.000107 & 0.943 & 0.523 & 0.0369 & 0.4559 & 0.6885 & 0.043 / 1.784 \\
Y2 & input & 1e-06 & 1554 & 0.061 & 0.219 & 0.158 & 0.0416 & 0.4849 & 0.9183 & 0.886 / 1.632 \\
Y2 & input & 0.0001 & 132 & 0.211 & 0.0143 & 0.0127 & 0.0845 & 0.6325 & 4.4719 & 1.516 / 1.580 \\
Y2 & input & 0.01 & 10 & 0.273 & 0.00115 & 0.00117 & 0.227 & 0.8581 & 26.9965 & 1.773 / 1.777 \\
Y2 & feature & 1e-08 & 2339 & 0.000906 & 0.451 & 0.44 & 0.00562 & 0.1431 & 0.1439 & 0.184 / 0.738 \\
Y2 & feature & 1e-06 & 564 & 0.0159 & 0.0893 & 0.0618 & 0.0143 & 0.2768 & 0.2696 & 0.662 / 0.899 \\
Y2 & feature & 0.0001 & 86 & 0.134 & 0.0118 & 0.00829 & 0.0512 & 0.5764 & 1.0187 & 1.275 / 1.326 \\
Y2 & feature & 0.01 & 9 & 0.27 & 0.000998 & 0.000909 & 0.188 & 0.8905 & 15.5377 & 1.738 / 1.741 \\
Y3 & input & 1e-08 & 3308 & 0.00797 & 0.7 & 0.412 & 0.0132 & 0.2382 & 0.3731 & 0.324 / 2.147 \\
Y3 & input & 1e-06 & 558 & 0.283 & 0.0644 & 0.0523 & 0.0192 & 0.2627 & 0.9135 & 1.656 / 1.951 \\
Y3 & input & 0.0001 & 45 & 0.455 & 0.00546 & 0.00501 & 0.0457 & 0.3432 & 5.4172 & 2.173 / 2.155 \\
Y3 & input & 0.01 & 5 & 0.554 & 0.000642 & 0.000643 & 0.152 & 0.4914 & 19.3748 & 2.715 / 2.702 \\
Y3 & feature & 1e-08 & 2371 & 0.00158 & 0.461 & 0.458 & 0.00569 & 0.0999 & 0.1197 & 0.221 / 0.903 \\
Y3 & feature & 1e-06 & 590 & 0.0255 & 0.094 & 0.063 & 0.0145 & 0.1650 & 0.1866 & 0.829 / 1.136 \\
Y3 & feature & 0.0001 & 90 & 0.267 & 0.0123 & 0.00842 & 0.0513 & 0.3145 & 0.6620 & 1.790 / 1.882 \\
Y3 & feature & 0.01 & 9 & 0.525 & 0.001 & 0.000885 & 0.189 & 0.4913 & 7.4936 & 2.612 / 2.624 \\
Y4 & input & 1e-08 & 3998 & 9.17e-07 & 0.999 & 0.473 & 0.269 & 2.1826 & 2.3706 & 0.007 / 12.844 \\
Y4 & input & 1e-06 & 3792 & 0.00736 & 0.9 & 0.437 & 0.27 & 2.1557 & 2.3495 & 0.629 / 12.779 \\
Y4 & input & 0.0001 & 1074 & 1.16 & 0.125 & 0.0912 & 0.319 & 1.9633 & 2.4907 & 7.986 / 11.573 \\
Y4 & input & 0.01 & 51 & 2.3 & 0.00312 & 0.00306 & 0.619 & 2.3401 & 8.9051 & 12.340 / 12.554 \\
Y4 & feature & 1e-08 & 3823 & 0.00807 & 0.917 & 0.975 & 0.031 & 2.4285 & 2.4310 & 0.626 / 14.459 \\
Y4 & feature & 1e-06 & 1652 & 0.794 & 0.272 & 0.24 & 0.0392 & 1.8004 & 1.8078 & 6.037 / 11.377 \\
Y4 & feature & 0.0001 & 229 & 1.83 & 0.0293 & 0.0257 & 0.103 & 1.8454 & 1.9424 & 9.636 / 10.497 \\
Y4 & feature & 0.01 & 19 & 2.39 & 0.0019 & 0.00181 & 0.323 & 2.4706 & 4.7750 & 12.907 / 12.992 \\
\end{longtable}}
\end{landscape}


\clearpage
\begin{table}[p]
\centering
% Generated by code/format_publication_tables.py; numerical cells preserved.
{\fontsize{9}{11}\selectfont
\setlength{\tabcolsep}{6pt}
\begin{tabular}{rrrrr}
\toprule
$\beta$ & \thead{average-family\\measured slope} & $\beta/(\beta+2)$ & \thead{uniform-family\\measured slope} & \thead{analytic slope\\same ladder} \\
\midrule
0.5 & 0.200 & 0.200 & 0.494 & 0.494 \\
1 & 0.333 & 0.333 & 0.973 & 0.973 \\
2 & 0.498 & 0.500 & 1.721 & 1.724 \\
3 & 0.597 & 0.600 & 1.963 & 1.963 \\
4 & 0.664 & 0.667 & 1.995 & 1.995 \\
\bottomrule
\end{tabular}}
\par\bigskip
{\fontsize{9}{11}\selectfont
\setlength{\tabcolsep}{8pt}
\begin{tabular}{rrrr}
\toprule
$\beta$ & \thead{uniform-family\\maximum relative gap} & \thead{profile-family\\maximum relative gap} & violations \\
\midrule
0.5 & 5.7e-04 & 4.4e-04 & 0 \\
1 & 1.7e-03 & 3.9e-04 & 0 \\
2 & 1.0e-02 & 3.1e-04 & 0 \\
3 & 1.4e-03 & 5.2e-04 & 0 \\
4 & 2.8e-04 & 2.8e-04 & 0 \\
\bottomrule
\end{tabular}}

\caption{Manufactured transfer families, seed zero, two million samples for each $\beta$. The
first panel compares finite-ladder slopes with the analytic curve on the same ladder. The
second reports the largest relative discrepancies between the Monte Carlo mean squared error and
the analytic value (uniform family) or the quadrature value (profile family), and the total
number of pointwise violations.}
\label{tab:v2-synthetic}
\end{table}
\clearpage
% Generated by code/format_publication_tables.py; numerical cells preserved.
\clearpage
\begin{landscape}
{\fontsize{9}{11}\selectfont
\setlength{\tabcolsep}{4pt}
\setlength{\LTcapwidth}{0.95\linewidth}
\begin{longtable}{lllrrrr}
\caption{Matched correction-coefficient benchmark \cite{pkan}. The standard split uses 35,000 / 7,500 / 7,500 states. OOD test states exceed the 0.85 aerosol or water-vapor quantile; remaining states are split 85/15. Cells average three seeds or initializations; published reference rows are not rerun. RMSE, MAE and SMAPE are pooled; $R^2$ is uniformly averaged across the three coefficients, reconstructed and checked from archived summary identities. Model keys: dkr, feature kernel; krr, input kernel; mlp, network; ens, ensemble; resid, residual correction; lowfi, the low-fidelity mean. The categorical inputs follow the benchmark protocol.}\label{tab:v2-pkanrtm}\\
\toprule
split & inputs & head & RMSE & MAE & $R^2$ & SMAPE [\%] \\
\midrule
\endfirsthead
\multicolumn{7}{l}{\tablename\ \thetable\ (continued)}\\
\toprule
split & inputs & head & RMSE & MAE & $R^2$ & SMAPE [\%] \\
\midrule
\endhead
\midrule
\multicolumn{7}{r}{Continued on next page}\\
\endfoot
\bottomrule
\endlastfoot
OOD & state only & dkr & 0.01408$\pm$0.00177 & 0.00709$\pm$0.00083 & 0.99226$\pm$0.00106 & 5.714$\pm$0.441 \\
OOD & state only & krr & 0.11283$\pm$0.00069 & 0.05430$\pm$0.00197 & 0.66968$\pm$0.01490 & 30.762$\pm$1.160 \\
OOD & state only & mlp & 0.00750$\pm$0.00044 & 0.00355$\pm$0.00036 & 0.99586$\pm$0.00012 & 4.500$\pm$0.249 \\
OOD & state only & mlp\_ens & 0.00624$\pm$0.00015 & 0.00258$\pm$0.00017 & 0.99667$\pm$0.00004 & 3.470$\pm$0.370 \\
OOD & state only & mlp\_resid & 0.00620$\pm$0.00015 & 0.00252$\pm$0.00015 & 0.99670$\pm$0.00005 & 3.372$\pm$0.308 \\
OOD & state only & ridge & 0.11944$\pm$0.00000 & 0.06195$\pm$0.00000 & 0.49009$\pm$0.00000 & 37.228$\pm$0.000 \\
OOD & with 6S inputs & dkr & 0.00735$\pm$0.00036 & 0.00336$\pm$0.00033 & 0.99597$\pm$0.00022 & 4.101$\pm$0.474 \\
OOD & with 6S inputs & krr & 0.02904$\pm$0.00015 & 0.01432$\pm$0.00005 & 0.97160$\pm$0.00020 & 11.238$\pm$0.033 \\
OOD & with 6S inputs & lowfi\_resid & 0.02901$\pm$0.00016 & 0.01430$\pm$0.00006 & 0.97172$\pm$0.00021 & 11.416$\pm$0.047 \\
OOD & with 6S inputs & mlp & 0.00741$\pm$0.00080 & 0.00347$\pm$0.00063 & 0.99602$\pm$0.00048 & 4.406$\pm$0.366 \\
OOD & with 6S inputs & mlp\_ens & 0.00611$\pm$0.00030 & 0.00246$\pm$0.00023 & 0.99682$\pm$0.00013 & 3.009$\pm$0.093 \\
OOD & with 6S inputs & mlp\_resid & 0.00596$\pm$0.00031 & 0.00232$\pm$0.00022 & 0.99687$\pm$0.00012 & 2.921$\pm$0.067 \\
OOD & with 6S inputs & ridge & 0.08048$\pm$0.00000 & 0.04396$\pm$0.00000 & 0.79274$\pm$0.00000 & 29.064$\pm$0.000 \\
standard & state only & dkr & 0.00874$\pm$0.00121 & 0.00387$\pm$0.00071 & 0.99379$\pm$0.00053 & 5.342$\pm$0.323 \\
standard & state only & krr & 0.11288$\pm$0.00001 & 0.05157$\pm$0.00002 & 0.68464$\pm$0.00028 & 31.178$\pm$0.032 \\
standard & state only & mlp & 0.00657$\pm$0.00026 & 0.00242$\pm$0.00030 & 0.99465$\pm$0.00007 & 5.136$\pm$0.550 \\
standard & state only & mlp\_ens & 0.00591$\pm$0.00007 & 0.00170$\pm$0.00005 & 0.99513$\pm$0.00010 & 3.873$\pm$0.132 \\
standard & state only & mlp\_resid & 0.00581$\pm$0.00004 & 0.00151$\pm$0.00004 & 0.99520$\pm$0.00010 & 3.543$\pm$0.066 \\
standard & state only & ridge & 0.11959$\pm$0.00000 & 0.06067$\pm$0.00000 & 0.48572$\pm$0.00000 & 41.647$\pm$0.000 \\
standard & with 6S inputs & dkr & 0.00627$\pm$0.00025 & 0.00218$\pm$0.00021 & 0.99502$\pm$0.00028 & 4.429$\pm$0.169 \\
standard & with 6S inputs & krr & 0.02514$\pm$0.00010 & 0.01155$\pm$0.00003 & 0.97447$\pm$0.00010 & 10.184$\pm$0.037 \\
standard & with 6S inputs & lowfi\_resid & 0.02513$\pm$0.00010 & 0.01155$\pm$0.00003 & 0.97452$\pm$0.00010 & 10.178$\pm$0.041 \\
standard & with 6S inputs & mlp & 0.00632$\pm$0.00018 & 0.00224$\pm$0.00017 & 0.99504$\pm$0.00014 & 4.496$\pm$0.160 \\
standard & with 6S inputs & mlp\_ens & 0.00587$\pm$0.00014 & 0.00163$\pm$0.00012 & 0.99531$\pm$0.00011 & 3.407$\pm$0.130 \\
standard & with 6S inputs & mlp\_resid & 0.00578$\pm$0.00012 & 0.00148$\pm$0.00008 & 0.99535$\pm$0.00010 & 3.185$\pm$0.097 \\
standard & with 6S inputs & ridge & 0.06981$\pm$0.00000 & 0.03631$\pm$0.00000 & 0.83813$\pm$0.00000 & 28.300$\pm$0.000 \\
standard & published & pKANrtm & 0.00619 & 0.00186 & 0.99472 & 3.912 \\
OOD & published & pKANrtm & 0.01022 & 0.00544 & 0.99418 & 5.399 \\
\end{longtable}}
\end{landscape}

% Generated by code/format_publication_tables.py; numerical cells preserved.
\clearpage
{\fontsize{9}{11}\selectfont
\setlength{\tabcolsep}{4pt}
\setlength{\LTcapwidth}{0.95\linewidth}
\begin{longtable}{llrr}
\caption{OCO-2 selection comparison, seeds 0--2 on each of three bands. Each record uses 2,000 validation and 2,000 test rows. Both reported columns are \emph{mean relative norms} in percent; the stack fitting objectives are row-relative \emph{squared} errors. The radiance quadratic uses the exact affine reconstruction Gram matrix. Model keys: dkr, feature kernel; mean, trained network; ridge, linear readout; raw/ard, input kernels. Flat, radx and wnum are the three training losses. Selection is per coordinate, with the radiance variant using greedy coordinate updates, not exhaustive global search.}\label{tab:v2-oco2}\\
\toprule
band & head / rule & reduced error [\%] & radiance error [\%] \\
\midrule
\endfirsthead
\multicolumn{4}{l}{\tablename\ \thetable\ (continued)}\\
\toprule
band & head / rule & reduced error [\%] & radiance error [\%] \\
\midrule
\endhead
\midrule
\multicolumn{4}{r}{Continued on next page}\\
\endfoot
\bottomrule
\endlastfoot
o2 & select\_radiance & 4.08$\pm$0.05 & 0.0283$\pm$0.0017 \\
o2 & select\_reduced & 4.09$\pm$0.06 & 0.0283$\pm$0.0017 \\
o2 & single dkr\_flat & 4.07$\pm$0.06 & 0.0757$\pm$0.0012 \\
o2 & single dkr\_radx & 20.14$\pm$0.07 & 0.0351$\pm$0.0006 \\
o2 & single dkr\_wnum & 21.70$\pm$0.60 & 0.0309$\pm$0.0017 \\
o2 & single kernel\_ard & 30.02$\pm$0.03 & 0.2706$\pm$0.1030 \\
o2 & single kernel\_raw & 40.53$\pm$0.09 & 0.2353$\pm$0.0047 \\
o2 & single mean\_flat & 4.40$\pm$0.08 & 0.1756$\pm$0.0111 \\
o2 & single mean\_radx & 58.72$\pm$0.60 & 0.0510$\pm$0.0049 \\
o2 & single mean\_wnum & 49.22$\pm$0.34 & 0.0437$\pm$0.0002 \\
o2 & single ridge\_flat & 4.67$\pm$0.09 & 0.2205$\pm$0.0124 \\
o2 & single ridge\_radx & 44.83$\pm$0.43 & 0.0538$\pm$0.0044 \\
o2 & single ridge\_wnum & 45.55$\pm$0.73 & 0.0527$\pm$0.0016 \\
o2 & stack\_radiance & 18.59$\pm$0.25 & 0.0264$\pm$0.0007 \\
o2 & stack\_reduced & 4.10$\pm$0.07 & 0.0291$\pm$0.0029 \\
sco2 & select\_radiance & 8.09$\pm$0.02 & 0.0598$\pm$0.0013 \\
sco2 & select\_reduced & 8.09$\pm$0.02 & 0.0598$\pm$0.0013 \\
sco2 & single dkr\_flat & 8.09$\pm$0.02 & 0.1083$\pm$0.0114 \\
sco2 & single dkr\_radx & 15.27$\pm$0.25 & 0.0780$\pm$0.0094 \\
sco2 & single dkr\_wnum & 15.48$\pm$0.31 & 0.0638$\pm$0.0022 \\
sco2 & single kernel\_ard & 28.72$\pm$0.02 & 0.6023$\pm$0.0097 \\
sco2 & single kernel\_raw & 41.01$\pm$0.09 & 0.5229$\pm$0.0081 \\
sco2 & single mean\_flat & 8.23$\pm$0.01 & 0.2107$\pm$0.0194 \\
sco2 & single mean\_radx & 46.08$\pm$0.91 & 0.0936$\pm$0.0082 \\
sco2 & single mean\_wnum & 36.64$\pm$0.38 & 0.0750$\pm$0.0020 \\
sco2 & single ridge\_flat & 8.34$\pm$0.02 & 0.2568$\pm$0.0244 \\
sco2 & single ridge\_radx & 41.14$\pm$1.78 & 0.0967$\pm$0.0081 \\
sco2 & single ridge\_wnum & 40.25$\pm$0.24 & 0.0887$\pm$0.0034 \\
sco2 & stack\_radiance & 8.87$\pm$0.15 & 0.0511$\pm$0.0016 \\
sco2 & stack\_reduced & 8.09$\pm$0.01 & 0.0536$\pm$0.0020 \\
wco2 & select\_radiance & 16.26$\pm$0.07 & 0.0510$\pm$0.0086 \\
wco2 & select\_reduced & 16.26$\pm$0.07 & 0.0510$\pm$0.0086 \\
wco2 & single dkr\_flat & 16.26$\pm$0.07 & 0.1135$\pm$0.0069 \\
wco2 & single dkr\_radx & 48.28$\pm$0.28 & 0.0475$\pm$0.0055 \\
wco2 & single dkr\_wnum & 48.60$\pm$0.38 & 0.0516$\pm$0.0086 \\
wco2 & single kernel\_ard & 52.45$\pm$0.07 & 0.2231$\pm$0.0014 \\
wco2 & single kernel\_raw & 58.63$\pm$0.20 & 0.4431$\pm$0.0063 \\
wco2 & single mean\_flat & 16.42$\pm$0.06 & 0.2254$\pm$0.0161 \\
wco2 & single mean\_radx & 67.92$\pm$0.46 & 0.0605$\pm$0.0075 \\
wco2 & single mean\_wnum & 62.39$\pm$0.42 & 0.0700$\pm$0.0087 \\
wco2 & single ridge\_flat & 16.49$\pm$0.07 & 0.2571$\pm$0.0234 \\
wco2 & single ridge\_radx & 61.03$\pm$0.31 & 0.0663$\pm$0.0080 \\
wco2 & single ridge\_wnum & 61.33$\pm$0.05 & 0.0778$\pm$0.0095 \\
wco2 & stack\_radiance & 37.43$\pm$0.20 & 0.0391$\pm$0.0026 \\
wco2 & stack\_reduced & 16.28$\pm$0.07 & 0.0404$\pm$0.0034 \\
\end{longtable}}

\clearpage


\section{Proofs}\label{app:proofs}

Numbers without the prefix S refer to the article.

\begin{proof}[Proof of Proposition~\ref{prop:stability}]
Subtract the two forward maps and use
$1/\widehat q-1/q=\rho e_s/(q\widehat q)$ to obtain
\eqref{eq:forward-error}. Since $u=\rho D$,
$\widehat D=D+e_t+u e_s-\widehat s e_a$ and
$\widehat\rho-\rho=(u-e_a-\rho\widehat D)/\widehat D$, which gives
\eqref{eq:inverse-error}. The triangle inequality yields
$|\widehat D-D|\le B$ and then \eqref{eq:inverse-bound}.
\end{proof}

\begin{proof}[Proof of Proposition~\ref{prop:identity}]
Subtract $h+S(f-h)$ from $f$ to obtain \eqref{eq:residual-exact}. On the design,
$I-K(K+\lambda I)^{-1}=\lambda(K+\lambda I)^{-1}$, which proves
\eqref{eq:design-residual} and its spectral form. Apply the assumed uniform inequality to
$g=f-h$ for \eqref{eq:inherit}.
\end{proof}

\begin{proof}[Proof of Proposition~\ref{prop:alignment}]
For a common operator the corrected error is $a-b$, so expansion of its squared norm gives
\eqref{eq:alignment}. The error of $p_t$ is $a-tb$; minimize the resulting quadratic on
$[0,1]$ to obtain \eqref{eq:shrink-mean}. For distinct operators,
$f-h-S_1(f-h)=(I-S_0)f-[(S_1-S_0)f+(I-S_1)h]$, after which the same expansion applies.
\end{proof}

\begin{proof}[Proof of Proposition~\ref{prop:sensitivity}]
(i) is the chain rule applied to $\widehat y=S_\lambda y+(I-S_\lambda)h_X(y)$, and the two examples are direct
substitutions. (ii) With the mean and kernel frozen the map is affine with linear part $S_\lambda$; the trace is the
sum of the eigenvalues $\mu_j/(\mu_j+\lambda)$, and the operator norm of a symmetric positive semidefinite matrix is
its largest eigenvalue. (iii) $\widehat y-f_X=-(I-S_\lambda)r+S_\lambda\varepsilon$; the cross term has zero mean,
and both matrices are diagonal in the eigenbasis of $K$. (iv) The correction at $x$ is $a_\lambda(x)^{\top}(y-h_X)$,
linear in $\varepsilon$ with coefficient vector $a_\lambda(x)$.
\end{proof}

\begin{proof}[Proof of Proposition~\ref{prop:selection}]
Let $\widehat{\mathcal R}_j$ and $\mathcal R_j$ be the empirical and population risks.
Hoeffding's inequality \cite{hoeffding} and a union bound give
$|\widehat{\mathcal R}_j-\mathcal R_j|\le
B\sqrt{\log(2M/\delta)/(2m)}$ simultaneously for all candidates with probability at least
$1-\delta$. Apply this once to the empirical minimizer and once to a population minimizer,
and use empirical optimality between the two inequalities.
\end{proof}

\begin{proof}[Proof of Theorem~\ref{thm:transfer}]
Write $q=1-\rho s$, $u=L-a=\rho t/q$, $D=t/q$, $\widehat D=\widehat t+\widehat s(L-\widehat a)$ and $z=\bar\rho$.
Suppose first $\widehat D>0$. The map $F(v)=\widehat Dv-(L-\widehat a)$ is increasing and affine, $z$ is the
projection of its root onto $[0,R]$ and $\rho\in[0,R]$, so $(z-\rho)F(z)\le0$: at an interior root $F(z)=0$; at
$z=R$ the root exceeds $R$, so $F(R)<0$ and $z-\rho\ge0$; at $z=0$ the root is negative, so $F(0)>0$ and
$z-\rho\le0$. Since $L-\widehat a=u-e_a$ and, by \eqref{eq:inverse-error},
$\widehat D=D+e_t+ue_s-\widehat se_a$ with $u=\rho D$,
\[
F(z)=D(z-\rho)+(1-z\widehat s)e_a+ze_t+zue_s ,
\]
hence $D(z-\rho)^2\le|z-\rho|\bigl((1-z\widehat s)|e_a|+z|e_t|+zu|e_s|\bigr)$ and, with $D=t/q$,
\[
|z-\rho|\le q(1-z\widehat s)\frac{|e_a|}{t}+qz\frac{|e_t|}{t}+z\rho\,|e_s| .
\]
Now $q(1-z\widehat s)$ is positive and at most $1$ under the hypotheses, which is all that is used here,
$qz\le R$ and $z\rho\le R^2$, so $|z-\rho|\le e_R/t$; and $|z-\rho|\le R$
because both lie in $[0,R]$. If $\widehat D\le0$, then $\widehat t\le\widehat s(e_a-u)\le\widehat s|e_a|$ and
$t=\widehat t-e_t\le|e_t|+\widehat s|e_a|$, so $Rt\le R|e_t|+R\widehat s|e_a|\le R|e_t|+|e_a|\le e_R$; the fallback
value and $\rho$ lie in $[0,R]$, so $|z-\rho|\le R=\min\{R,e_R/t\}$. For optimality take $s=\widehat s=0$,
$\widehat t=t$, $\rho=r_0$ and $\widehat a=a+ct$ with $0<c<r_0<R$: then $|\bar\rho-\rho|=c=e_R/t$ exactly, with a
positive denominator and no projection active. This proves (i). For (ii) integrate (i), split at $t=\tau$ and bound
$e_R^2/t^2$ by $e_R^2/\tau^2$ on $\{t>\tau\}$. For (iii), $\varepsilon_R=0$ gives zero error by (i). Otherwise take $\tau=\varepsilon_R^{2/(\beta+2)}\le u_0$; then
$R^2A\tau^\beta+\varepsilon_R^2/\tau^2=(R^2A+1)\varepsilon_R^{2\beta/(\beta+2)}$. (iv) is (i) with $t\ge t_0$.
\end{proof}

\begin{proof}[Proof of Proposition~\ref{prop:attained}]
$L=r_0t$ and $\widehat\rho=(r_0t-\widehat a)/t$; $e_R=|\widehat a|=ct\,\1\{t\le\tau\}$; with density
$\beta u^{\beta-1}$ on $(0,1)$, $\E_P[t^2\1\{t\le\tau\}]=\beta\tau^{\beta+2}/(\beta+2)$; the retrieval error equals
$c$ exactly on $\{t\le\tau\}$, an event of probability $\tau^\beta$. Eliminate $\tau$.
\end{proof}

\begin{proof}[Proof of Theorem~\ref{thm:profile}]
If $R=0$ the statement is immediate; assume $R>0$. Conditional on $t$ the squared error is at most both $R^2$ and $e_R^2/t^2$ by Theorem~\ref{thm:transfer}(i), so
$\E[|\bar\rho-\rho|^2\mid t]\le\min\{R^2,\E[e_R^2\mid t]/t^2\}\le\min\{R^2,\sigma^2t^{2\nu-2}\}$; integrate. With
$t_*=(\sigma/R)^{1/(1-\nu)}$,
$\E_P\min\{R^2,\sigma^2t^{2\nu-2}\}=R^2F_t(t_*)+\sigma^2\E_P[t^{2\nu-2}\1\{t>t_*\}]$; the first term is at most
$R^2At_*^\beta=O(\sigma^{\beta/(1-\nu)})$ once $t_*\le u_0$, and
$\E_P[t^{2\nu-2}\1\{t>t_*\}]=\int_0^{t_*^{2\nu-2}}\Prob(t_*<t<v^{-1/(2-2\nu)})\,dv
\le u_0^{2\nu-2}+A\int_{u_0^{2\nu-2}}^{t_*^{2\nu-2}}v^{-\beta/(2-2\nu)}\,dv$, which is $O(1)$ for
$\beta>2(1-\nu)$, $O(\log(1/\sigma))$ for $\beta=2(1-\nu)$ and $O(t_*^{\beta-2(1-\nu)})$ for $\beta<2(1-\nu)$. At
$\nu=1$ the bound is $\sigma^2$ before any tail enters. In the model,
$c^{-2}\E_P|\bar\rho-\rho|^2=\E_P\min\{1,(\sigma/c)^2t^{2\nu-2}\}$, the same integral with equality.
\end{proof}

\begin{proof}[Proof of Proposition~\ref{prop:stackweights}]
(i) With $a=s=0$ and exact $t$ and $s$, the stack's path-radiance error is
$e(w)=((1-w)b,\,wc\delta)$, and the radiance at $\rho$ is predicted with the same error, since
$q=\widehat q=1$; this gives $F$. The constrained retrieval evaluates
$(L-\widehat a)/t=\rho-e_i/t_i$ at entry $i$, with denominator $t_i>0$, and
$|e_i|/t_i\le\max\{b,c\}\le\min\{\rho,R-\rho\}$ keeps it in $[0,R]$, so the projection is inactive
and $\bar\rho-\rho=-((1-w)b,\,wc)$; this gives $E$. Both are convex quadratics in $w$, and setting
their derivatives to zero gives $w_f$ and $w_r$. Substituting,
\[
E(w_f)=\frac{b^2c^2\,(c^2\delta^4+b^2)}{2(b^2+c^2\delta^2)^2},\qquad
E(w_r)=\frac{b^2c^2}{2(b^2+c^2)},
\]
whose ratio is the one stated. It tends to $(b^2+c^2)/b^2$ as $\delta\to0$; at $c=b/\delta$ the
two factors of the numerator are $b^2(1+\delta^2)$ and $b^2(1+\delta^2)/\delta^2$ and the
denominator is $4b^4$. Since $w_f$ minimizes $F$, $F(w_f)\le F(w_r)$. Finally, with
$\widetilde e_t=\widetilde e_s=0$ and $\tau\le\delta\le1$, the weights $(t_i\vee\tau)^{-2}$ in
\eqref{eq:Jtau} are $1$ and $\delta^{-2}$, so $J_\tau(w)=\tfrac12[(1-w)^2b^2+w^2c^2\delta^2/\delta^2]=E(w)$.

(ii) With $s=\widehat s=0$ and $\widehat t=t$, the predicted radiance at $\rho$ is
$\widehat a+\rho t$, so its error is the path-radiance error $r_w$ of the combination, and the
retrieval evaluates $(L-\widehat a)/t=\rho-r_w/t$, which lies in $[0,R]$ by hypothesis. Hence
$F(w)=\E_Pr_w^2$ and $E(w)=\E_P[r_w^2/t^2]$. Since the weights sum to one, $r_w$ lies in
$\mathcal L$, so $\lambda_{\min}F(w)\le E(w)\le\lambda_{\max}F(w)$ for every $w$ (both sides vanish
when $\E_Pr_w^2=0$), and
\[
E(w_f)\le\lambda_{\max}F(w_f)\le\lambda_{\max}F(w_r)\le\frac{\lambda_{\max}}{\lambda_{\min}}\,E(w_r).
\]
If $t_{\min}\le t\le t_{\max}$ almost surely, then
$r^2/t_{\max}^2\le r^2/t^2\le r^2/t_{\min}^2$ pointwise, so $\lambda_{\max}\le t_{\min}^{-2}$ and
$\lambda_{\min}\ge t_{\max}^{-2}$. In the example of (i), $\mathcal L=\mathbb R^2$ and the ratio
$\tfrac12(r_1^2+r_2^2/\delta^2)/\tfrac12(r_1^2+r_2^2)$ runs from $1$ to $\delta^{-2}$, while
$E(w_f)/E(w_r)$ at $c=b/\delta$ is at least $\delta^{-2}/4$.
\end{proof}

\begin{proof}[Proof of Proposition~\ref{prop:perturb}]
$\widetilde H=H^{1/2}(I+B)H^{1/2}$ gives $\widetilde H^{-1}=H^{-1/2}QH^{-1/2}$ and $Q=I-QB$. Hence
$\widetilde k_x^{\top}\widetilde H^{-1}(Y-\widetilde M)-k_x^{\top}H^{-1}(Y-M)
=\widetilde g^{\top}Q(W-U)-g^{\top}W
=\widetilde g^{\top}W-\widetilde g^{\top}QBW-\widetilde g^{\top}QU-g^{\top}W
=v^{\top}W-(Q\widetilde g)^{\top}(BW+U)$, which is (i) since $Q$ is symmetric; (ii) bounds its two terms. For (iii)
write the same difference as $v^{\top}QW+g^{\top}(Q-I)W-\widetilde g^{\top}QU$ and use
$\|Q\|_{\mathrm{op}}=\gamma_Q$ and $\|Q-I\|_{\mathrm{op}}=\|QB\|_{\mathrm{op}}=\eta_Q$. For (iv), $v=0$, $U=0$ and
$B=sH^{-1}$ has eigenvalues $s/\lambda_i(H)\ge0$, so $\eta_Q=\max_i(s/\lambda_i)/(1+s/\lambda_i)=s/(\lambda_{\min}(H)+s)$.
\end{proof}



\section*{Acknowledgements and funding}

The research presented in this paper was supported by the European Research Council (ERC) under
the European Union's Horizon Europe research and innovation programme (grant agreement
No.~101041711), by the Simons Foundation as part of the Collaboration on the Mathematical and
Scientific Foundations of Deep Learning, by Heights Labs, by Convex Nexus Capital, by the Israel
Science Foundation (grant number 2258/19), by the Israel Science Foundation (ISF Grant
4101/25), and by Google Cloud through its education credits program (credit 530550103).

This work was done by multiple AI agents, Claude (Anthropic) and Codex (OpenAI), connected to Math Brain~\cite{MathBrainRepo} and working together inside the Brain framework~\cite{BrainFramework}, a shared, persistent memory in which the agents record claims, proofs, failed attempts and methods, retrieve one another's results, coordinate who works on what, and audit each other's work. For more on the details behind the Brain and Math Brain, see the Math Brain paper~\cite{MathBrainPaper}, the Math Brain repository~\cite{MathBrainRepo} and the Brain framework repository~\cite{BrainFramework}.

\section*{Declaration of competing interest}

The author holds equity in Heights Labs and in Convex Nexus Capital and has a management role at
Convex Nexus Capital; both entities supported this work and are named above. These entities and
the public funders had no role in the study design, analysis, manuscript preparation, or
decision to publish.

\section*{AI-assisted contributions}

Claude Code (Anthropic) and Codex (OpenAI) implemented the data pipeline and the model
comparisons, ran the experiments, produced the figures and tables from the run records, and
drafted the manuscript. ChatGPT (OpenAI) developed the residual-error and inversion-stability
arguments, checked the identities and wrote the conditioning diagnostics and their tests. The
author supplied the data, set the questions, and reviewed the results, the proofs and the text.

\section*{Data and code availability}

The code, the run records and this supplement are in the repository
\url{https://github.com/yspennstate/emit-oco2-kernel-emulators}. The paper is part of the Math Brain project, whose public repository is \url{https://github.com/yspennstate/math-brain-public}~\cite{MathBrainRepo}.

\begin{thebibliography}{99}

\bibitem{lamminpaa}
O.~Lamminp\"a\"a, J.~Susiluoto, J.~Hobbs, J.~McDuffie, A.~Braverman, and H.~Owhadi.
\newblock Forward model emulator for atmospheric radiative transfer using Gaussian processes and
cross validation.
\newblock \emph{Atmospheric Measurement Techniques}, 18:673--694, 2025.

\bibitem{pkan}
M.~A.~A. Mazid and N.~Rishe.
\newblock Multi-fidelity emulation of atmospheric correction coefficients with physics-guided
Kolmogorov--Arnold networks.
\newblock \emph{Remote Sensing}, 18(11):1826, 2026.

\bibitem{rw}
C.~E. Rasmussen and C.~K.~I. Williams.
\newblock \emph{Gaussian Processes for Machine Learning}.
\newblock MIT Press, 2006.

\bibitem{nmkc}
Y.~Shmalo.
\newblock Neural means and kernel corrections for operator learning.
\newblock Preprint, ResearchGate, 2026. Code and per-run summaries at
\url{https://github.com/yspennstate/neural-means-kernel-corrections}.

\bibitem{vermote}
E.~F. Vermote, D.~Tanr\'e, J.~L. Deuz\'e, M.~Herman, and J.-J. Morcrette.
\newblock Second simulation of the satellite signal in the solar spectrum, 6S: an overview.
\newblock \emph{IEEE Transactions on Geoscience and Remote Sensing}, 35(3):675--686, 1997.

\bibitem{mayer}
B.~Mayer and A.~Kylling.
\newblock Technical note: the {libRadtran} software package for radiative transfer
calculations --- description and examples of use.
\newblock \emph{Atmospheric Chemistry and Physics}, 5:1855--1877, 2005.

\bibitem{owhadiyoo}
H.~Owhadi and G.~R. Yoo.
\newblock Kernel flows: from learning kernels from data into the abyss.
\newblock \emph{Journal of Computational Physics}, 389:22--47, 2019.

\bibitem{nystrom}
C.~K.~I. Williams and M.~Seeger.
\newblock Using the {Nystr\"om} method to speed up kernel machines.
\newblock In \emph{Advances in Neural Information Processing Systems 13}, pages 682--688, 2001.

\bibitem{wilsonadams}
A.~G. Wilson and R.~P. Adams.
\newblock {Gaussian} process kernels for pattern discovery and extrapolation.
\newblock In \emph{Proceedings of the 30th International Conference on Machine Learning},
pages 1067--1075, 2013.

\bibitem{additivegp}
D.~Duvenaud, H.~Nickisch, and C.~E. Rasmussen.
\newblock Additive {Gaussian} processes.
\newblock In \emph{Advances in Neural Information Processing Systems 24}, 2011.

\bibitem{stonecv}
M.~Stone.
\newblock Cross-validatory choice and assessment of statistical predictions.
\newblock \emph{Journal of the Royal Statistical Society: Series B}, 36(2):111--147, 1974.

\bibitem{gcv}
G.~H. Golub, M.~Heath, and G.~Wahba.
\newblock Generalized cross-validation as a method for choosing a good ridge parameter.
\newblock \emph{Technometrics}, 21(2):215--223, 1979.

\bibitem{arlot}
S.~Arlot and A.~Celisse.
\newblock A survey of cross-validation procedures for model selection.
\newblock \emph{Statistics Surveys}, 4:40--79, 2010.

\bibitem{adamw}
I.~Loshchilov and F.~Hutter.
\newblock Decoupled weight decay regularization.
\newblock In \emph{International Conference on Learning Representations}, 2019.

\bibitem{prechelt}
L.~Prechelt.
\newblock Early stopping --- but when?
\newblock In \emph{Neural Networks: Tricks of the Trade}, volume 1524 of \emph{Lecture Notes
in Computer Science}, pages 55--69. Springer, 1998.

\bibitem{leenngp}
J.~Lee, Y.~Bahri, R.~Novak, S.~S. Schoenholz, J.~Pennington, and J.~Sohl-Dickstein.
\newblock Deep neural networks as {Gaussian} processes.
\newblock In \emph{International Conference on Learning Representations}, 2018.

\bibitem{chosaul}
Y.~Cho and L.~K. Saul.
\newblock Kernel methods for deep learning.
\newblock In \emph{Advances in Neural Information Processing Systems 22}, pages 342--350, 2009.

\bibitem{ntk}
A.~Jacot, F.~Gabriel, and C.~Hongler.
\newblock Neural tangent kernel: convergence and generalization in neural networks.
\newblock In \emph{Advances in Neural Information Processing Systems 31}, 2018.

\bibitem{rfm}
A.~Radhakrishnan, D.~Beaglehole, P.~Pandit, and M.~Belkin.
\newblock Mechanism for feature learning in neural networks and backpropagation-free machine
learning models.
\newblock \emph{Science}, 383(6690):1461--1467, 2024.

\bibitem{dehoop}
M.~V. de~Hoop, D.~Z. Huang, E.~Qian, and A.~M. Stuart.
\newblock The cost-accuracy trade-off in operator learning with neural networks.
\newblock \emph{Journal of Machine Learning}, 1(3):299--341, 2022.

\bibitem{pdebench}
M.~Takamoto, T.~Praditia, R.~Leiteritz, D.~MacKinlay, F.~Alesiani, D.~Pfl\"uger, and
M.~Niepert.
\newblock {PDEBench}: an extensive benchmark for scientific machine learning.
\newblock In \emph{Advances in Neural Information Processing Systems 35}, 2022.

\bibitem{thewell}
R.~Ohana, M.~McCabe, L.~Meyer, R.~Morel, F.~Agocs, M.~Beneitez, M.~Berger, B.~Burkhart,
S.~Dalziel, D.~Fielding, et~al.
\newblock The Well: a large-scale collection of diverse physics simulations for machine
learning.
\newblock In \emph{Advances in Neural Information Processing Systems 37}, pages
44989--45037, 2024.

\bibitem{climsim}
S.~Yu, W.~M. Hannah, L.~Peng, J.~Lin, M.~A. Bhouri, R.~Gupta, B.~L\"utjens, J.~C. Will,
G.~Behrens, et~al.
\newblock {ClimSim}: a large multi-scale dataset for hybrid physics-{ML} climate emulation.
\newblock In \emph{Advances in Neural Information Processing Systems 36}, 2023.

\bibitem{golubvanloan}
G.~H. Golub and C.~F. Van~Loan.
\newblock \emph{Matrix Computations}.
\newblock Johns Hopkins University Press, fourth edition, 2013.

\bibitem{hoeffding}
W.~Hoeffding.
\newblock Probability inequalities for sums of bounded random variables.
\newblock \emph{Journal of the American Statistical Association}, 58(301):13--30, 1963.

\bibitem{stewartsun}
G.~W. Stewart and J.-G. Sun.
\newblock \emph{Matrix Perturbation Theory}.
\newblock Academic Press, 1990.

\bibitem{higham}
N.~J. Higham.
\newblock \emph{Accuracy and Stability of Numerical Algorithms}.
\newblock SIAM, second edition, 2002.

\bibitem{MathBrainRepo} Y.~Shmalo, \emph{Math Brain: public snapshot, code and tests}, 2026, \url{https://github.com/yspennstate/math-brain-public}.

\bibitem{MathBrainPaper} Y.~Shmalo, \emph{Math Brain: shared, auditable memory for multi-agent mathematical research}, October 2026, \url{https://github.com/yspennstate/math-brain-public/tree/main/papers/math-brain-2026-10-24}.

\bibitem{BrainFramework} Y.~Shmalo, \emph{Brain Framework}, public software repository, 2026, \url{https://github.com/yspennstate/brain-framework-public}.

\end{thebibliography}


\end{document}
